    \documentclass[../main.tex]{subfiles}
\begin{document}
%==========================================TIÊU ĐỀ TẬP========================================
\begin{table}[h]
  \begin{adjustwidth}{-1cm}{}
    \begin{tabular}{>{\centering\arraybackslash}p{6cm}|>{\raggedright\arraybackslash}p{12.5cm}}
      \multirow{5}{6cm}
      % Cột bên trái: ÉPISODE và số tập
      {\\[-40pt]
        \flaregothic\fontsize{20pt}{20pt}\selectfont \centering
        {\contour{errorfive}{\textcolor{black}{ÉPISODE}}}\color{black}\\
        \burbank\fontsize{72pt}{72pt}\selectfont
        \includegraphics[height=3.12359cm]{../../../../Image/Book?/number/num31.png}
        \\[18pt]
        % \flaregothic\fontsize{14pt}{14pt}\selectfont
        % \color{epattr}BIENVENUE, \uline{\color{Banpire} Mel} !
        % \flaregothic\fontsize{18pt}{18pt}\selectfont
        % \color{epattr}ÉPISODE PERDU
        \color{black}
      }
       &
      % Dòng thứ nhất: tên tập tiếng Nhật
      {\fotskip\fontsize{18pt}{18pt}\selectfont \ruby{過}{か}\ruby{去}{こ}の\ruby{残}{ざん}\ruby{響}{きょう}}
      \\[6pt]
      % Dòng thứ hai: tên tập tiếng Anh   
       & {\cafeteria\fontsize{18pt}{18pt}\selectfont Echoes of the Past}                                                                        \\[4pt]
      % Dòng thứ ba: tên tập tiếng Việt
       & {\vnmsans\fontsize{18pt}{18pt}\selectfont Tiếng vọng quá khứ}                                                                          \\[8pt]
      % Dòng thứ tư: nhân vật xuất hiện
       & {\caviar{\fontsize{14pt}{14pt}\selectfont Nhân vật: \emp{sora} \emp{miko} \emp{azki} \emp{shion} \emp{choco} \emp{marine} \emp{nene}}}
      % \\[6pt]
      % Dòng cuối cùng: Thời gian diễn ra của tập (thường sẽ là ngày phát hành)
      % & {\continuum\fontsize{16pt}{16pt}\selectfont Ngày phát hành: 11/06/2022}\\[8pt]
    \end{tabular}
  \end{adjustwidth}
\end{table}
%===========================================NỘI DUNG==========================================
\color{black}

\pov{Hikawa Kuon}

Vài ngày sau hội thao, mọi thứ dần quay trở lại quỹ đạo vốn có của một lớp học bình thường. Tiếng chuông reo vào giờ, giày thể thao bám mùi cỏ khô trên hành lang, bảng tin trước cổng dán lịch thi học kỳ, bảng lớp 1-C cập nhật ca trực nhật mới, và điểm danh vẫn vang lên đều đặn như nhịp tim của ngôi trường. Những học sinh của lớp vẫn cười giòn, giáo viên vẫn xoa phấn lên bảng, bầu không khí nhộn nhịp trong căng-tin mỗi lần chuông báo giờ nghỉ giải lao\dots\ Một mặt nạ bình yên được tái lập bằng những chi tiết nhỏ nhặt nhất, khiến bất kỳ khách qua đường nào cũng tin rằng chưa từng có cuộc chiến tí hon nào xảy ra.

Nhưng dĩ nhiên, tất cả chỉ là lớp sơn mới phủ lên vết nứt.

Tôi vẫn ngồi chỗ cũ, cửa sổ hắt ánh nắng lên bàn. Sân trường trống vắng, không còn bóng người chạy đua, chỉ còn lá phong đỏ lác đác rơi. Giảng đường 1-C trở lại nhịp thở quen: tiếng phấn gõ bảng, giấy xào xạc, lời thủ thỉ về bài kiểm tra sắp tới.

Một cuộc sống học sinh ``bình thường''. Nhưng tôi biết, đâu đó dưới lớp vỏ ấy, mạch máu đã đổi dòng.

``Cậu lại đang nghĩ vẩn vơ gì thế?'' một giọng nói vang lên từ bên cạnh. Kaigou-chan khoanh tay, hơi nghiêng đầu nhìn tôi, ánh mắt như muốn đọc hết suy nghĩ trong đầu tôi.

``Tôi đang nghĩ xem hôm nay cô có định nghe giảng không thôi.''

``Điêu.'' Cô thở dài sau đó, nhưng không nói thêm gì.

Trước mặt, Suzuki xoay ghế, đôi mắt hai màu giương lên với sự lo ngại: ``Onii-sama\dots\ anh cũng cảm thấy đúng không?''

Tôi gật, như thể gật cho chính mình.

Em ấy thả giọng xuống thấp hơn cả tiếng ve: ``Dạo gần đây, em thấy\dots\ chúng không còn lang thang vô định nữa.''

Đúng vậy, những ``đốm sáng'' - mảnh linh hồn trẻ con - giờ tuân theo một quy luật nhất định, chứ không còn xuất hiện đầy ngẫu nhiên như trước nữa. Tôi đã để ý trong những ngày vừa qua, khi buổi sáng đi học, khi thì giờ nghỉ, và khi tôi làm việc trong cửa hàng tiện lợi.

Vào buổi sáng, khuôn viên trường sạch bong, chỉ còn sương và tiếng chim.

Sau đó, tầm mười một giờ, chúng lảng vảng hành lang cũ kỹ phía sau nhà thể chất, như học sinh trốn tiết.

Lúc chiều muộn, mật độ tăng vọt quanh đền Aogami, đèn lồng chưa kịp thắp đã bật sáng bởi chính những đốm ấy.

Và khi hoàng hôn nhường chỗ cho màn đêm, chúng bắt đầu diễu hành về phía rừng xung quanh ngôi đền Aogami.

``\textit{Như có ai thổi sáo gọi đàn,}'' tôi lẩm bẩm.

Suzuki siết chặt tay mình, gân xanh nổi lên: ``Vâng\dots\ Em cũng thấy vậy.''

Chúng không còn lẩn khuất như trước. Những đốm sáng ấy giờ di chuyển thành dòng, lặng lẽ rời phố, băng qua sân trường, rồi hướng thẳng về triền núi phía sau đền, nơi bóng tối như đang há miệng chờ nuốt chửng ánh sáng.

``Mình cũng thấy chúng tối qua này,'' Shino-san ở bàn bên khẽ nói, vừa lật cuốn sổ nhỏ ra, ``mình đang đếm thử xem một đêm có bao nhiêu đốm, nhưng chúng bay nhanh quá, mình đếm lạc cả\dots''

Sakura-san ử đằng sau khẽ cất tiếng, giọng thấp hơn cả tiếng lá rơi: ``Tôi đã thử ghi lại quỹ đạo của chúng.'' Cô lấy trong cặp ra tờ giấy kẻ gấp bốn. Mở ra, bản phác thị trấn Aogami hiện lên, phủ đầy chấm nhỏ li ti. Mỗi chấm là một lần chúng xuất hiện, được ghi lại bằng bút chì khắt khe.

Nhìn kỹ lại, chẳng cái nào rơi rụng tình cờ. Chúng xếp thành những cung tròn mờ, như bản đồ nhiệt của một dòng chảy ngầm, tất cả đang xoáy về phía rừng thông đằng xa. Đúng như tôi dự đoán lúc nãy.

Không khí chợt nặng. Suzuki nắm chặt viền áo, im như thóc.

*BÍP BÍP*

Tiếng động phát ra từ cổ tay phải của tôi. Chiếc đồng hồ của tôi bất chợt lóe sáng, như là một lần hiếm hoi được dịp thể hiện sau những ngày tháng im lìm vì không muốn bị tiết lộ thân phận.

Giọng nói trầm ấm, rõ ràng, vang lên từ dây da: ``\eng{Kuon-user, tôi đã xác minh. Các linh hồn trẻ con đang di chuyển theo một vector duy nhất: rừng phía sau đền Aogami. Tần suất tăng 340\% trong ba ngày qua. Điểm hội tụ nằm sâu bên trong, cách ngôi đền khoảng 400 mét.}''

Kaigou-chan chau mày: ``\eng{Ông có thể định lượng được ư?}''

``\eng{Tôi vừa làm xong. Phân bố không còn ngẫu nhiên; chúng tuân theo đường xoắn ốc Fibonacci, bán kính thu hẹp 12 mét mỗi giờ. Nếu duy trì tốc độ này, 21:47 tối nay chúng sẽ biến mất khỏi vùng quét của tôi.}''

Shino-san run nhẹ: ``\eng{Biến mất\dots\ là sao?}''

``\eng{Hoặc bị nuốt chửng, hoặc hợp nhất thành một thực thể duy nhất. Tôi chưa thể xác định được. Nhưng tôi nghĩ, tốt nhất các cậu nên tìm hiểu thêm, càng sớm càng tốt.}''

``\eng{Hiểu rồi.}''

Mặt kính lại trở về vẻ trơ tuồng, nhưng tôi biết: bản đồ Kaigou vẽ chỉ là bề nổi, còn rừng kia đã sẵn sàng lên đèn, và chúng tôi cần phải tìm ra mục đích thực sự của chúng là gì.

``\textbf{Kuon-kun!!}''

Tiếng gọi vang lên, nhẹ nhưng đủ kéo tôi khỏi mớ suy tưởng đang cuộn. Mari-san đứng cuối lớp, ôm cuốn sổ vẽ vào lòng; nụ cười vẫn tươi, nhưng khóe mắt hơi giật, như thể có con sóng nhỏ vừa đập vào bờ ý thức cô.

``Cậu rảnh chứ? Cậu xem thử cái này đi.''

Tôi gật, đẩy ghế ra. Mỗi bước tới chỗ cô như đi trên mặt nước.

Cô nàng Họa sĩsĩ lật mở cuốn sổ. Trang đầu tiên, cô vẽ một cánh rừng. Không phải rừng mà tôi hay bất cứ người nào mừng tượng một cách thông thường. Từng vân gỗ, từng mảng sáng bị cây cốa cắt xé, kể cả mùi đất ẩmẩm cũng như toát ra khỏi nét bút chì. Nó chân thực tới độ tôi ngửi thấy hơi đất dù chỉ đang nhìn.

``Tôi không nhớ là đã từng tới đó,'' cô ấy nói, giọng nhỏ lại, ``nhưng chẳng sao\dots\ tay cứ tự vẽ.''

Tôi chưa kịp đáp, cô ấy đã lật tiếp. Đó là một đôi mắt. Xanh lá sâu thẳm, đục ngầu như hồ nước tù đọng mùa đông. Đôi mắt không vẽ con ngươi, chỉ toàn màu, nhưng tôi vẫn cảm được nó đang dõi theo mình.

\dots Đó chẳng phải là đôi mắt của cô bé hồn ma đó sao?

Suzuki bên cạnh rụt vai, hơi thở gấp thành từng đợt, đôi mắt em ấy hơi co lại.

``Kỳ lạ đúng không? Tôi cũng không hiểu tại sao mình lại cứ vẽ đi vẽ lại cái này. Mỗi lần mở sổ ra là lại thấy nó ở đó rồi. Như thể\dots'' cô ngưng lại một chút, rồi tiếp tục, ``\dots thế giới này nhét ý mình vào ngòi bút.''

Tôi thấy gáy mình tê dại. Không chỉ Mari-san. Tôi cũng từng thấy khu rừng ấy, trong giấc mộng hôm trước, hay trong một ký ức bị xóa nhòa, tôi không còn biết nữa. Một sự thân quen bám lấy tôi như mảng sương giăng trên tán cây.

Hình ảnh cũ bỗng tràn về, nhòe nhoẹt như bức tranh bị xé xác: căn phòng mỹ thuật ngập mùi dầu hoả, từng nét vẽ trên giá đứng run rẩy dưới ánh đèn neon, và nụ cười của Mari-san khi máu từ khóe môi rỉ ra, thỏa mãn đến rợn người. Tôi vẫn nghe thấy tiếng Suzuki (hồi đó vẫn là Reiko) la hét hoảng loạn, tiếng ai đó thì thầm ``kiệt tác'' giữa vũng đỏ loang dưới chân giá vẽ, và bức tranh đáng nguyền rủa đó.

Tôi gập cuốn sổ lại, mạnh đến mức lòng bàn tay đau nhói. Có thứ gì đó trong những trang giấy vẫn thở, vẫn giật giật như muốn thoát ra. Tôi không cho phép. Không đời nào để cơn ác mộng ấy tái diễn, không đời nào để Mari-san trở thành con mồi của chính nghệ thuật cô yêu thương.

``Cậu nên tạm nghỉ vẽ một thời gian đi,'' giọng tôi khàn đặc, không cho phép cãi lại.

``Hả? Nhưng mà--'' Cô ôm sổ vẽ vào lòng, như thể tôi vừa đòi giật đi đứa con tinh thần.

``Nghe tôi đi.'' Tôi cúi xuống, ánh mắt ghim chặt vào đôi mắt cô, cố gắng truyền hết nỗi sợ của mình qua khoảng cách vài tấc. ``Bức tranh mà cậu vẽ \emph{không bình thường} đâu.'' Tôi không nói gì thêm, nhưng cũng đủ để khiến cô nàng tóc nâu-đỏ chùn bước.

Không khí giữa chúng tôi như dày đặc lại. Mari-san mím môi, lông mày nhíu chặt, rồi buông một tiếng thở dài gần như không nghe thấy. Cô gật đầu, chậm rãi, như thể từng chuyển động đều đánh cắp một phần linh hồn cô.

Tôi thấy tay cô run nhẹ khi vuốt lại bìa sổ. Có lẽ cô cũng vừa nhìn thấy bóng dáng của quá khứ bản thân lẩn khuất trong trang giấy trắng.

Mari-san ôm sổ vẽ trở về chỗ ngồi, bóng dáng cô nhỏ lại trong ánh đèn huỳnh quang. Tôi quay lại bàn, nhưng cảm giác nặng nề không buông tha; như thể có một sợi dây vô hình từ trang giấy vẫn quấn quanh cổ tay tôi, kéo tôi về phía khu rừng đang chờ đợi phía xa.

``Mình thấy bức vẽ đó cũng có chút hơi hướng linh hồn,'' Shino-san khẽ nói, ngón tay gõ nhịp lên cuốn sổ nhỏ, ``cậu có cảm thấy như nó đang \emph{thở} không?''

Sakura-san đáp lại: ``Tôi đoán năng lượng bám quanh trang giấy khá mạnh. Nếu cậu cần, tôi có thể thử định lượng tần số dao động.''

``Cảm ơn cậu,'' tôi gật đầu, ``nhưng đừng để nó kéo cậu vào.''

Cô nàng tóc nâu mỉm cười an ủi: ``Yên tâm, mình chỉ đếm những đốm sáng thôi mà.''

Cô nàng Vu nữ khẽ nhíu mày: ``Cậu nên cất cuốn sổ vào hộp chì, cách ly hoàn toàn. Tôi sẽ chuẩn bị lá bùa bảo vệ cho Mari-sansan nếu cậu ấy cần.''

``Mình sẽ giúp cậu ấy,'' Shino-san nói tiếp, ``mình có bút chì khắc bùa, chỉ cần vài nét là xong.''

``Đừng quên ghi lại dữ liệu,'' Sakura-sạn nhắc, ``nếu có gì bất thường, cậu báo ngay cho tôi.''

Hai cô gái gật đầu, ánh mắt đầy quyết tâm. Tôi biết chúng tôi không chỉ đối mặt với bức tranh, mà còn với cả khu rừng đang thức giấc phía sau những trang giấy.

\ct

Tiếng chuông vừa dứt, lớp học đã rộ lên những tiếng xì xầm. Không ai nhắc đến bức vẽ, cũng chẳng ai dám. Thay vào đó, họ nói về thứ vẫn âm ỉ trong gió đêm.

Azuki-san vỗ nhẹ lên mặt bàn, kéo ghế sát lại, đôi mắt cô sáng lên vẻ tò mò vừa hồn nhiên vừa man dại: ``Mình kể nghe nè, Kuon-kun! Mình vẫn còn vương vấn về lời đồn đại mà Honoka-san đã nói hôm trước!''

Tôi hơi nhướng mày. Nhưng tôi không nói gì thêm, tiếp tục để cô ấy kể. Cô khựng lại một nhịp, như thể cân nhắc xem có nên thêm chi tiết không, rồi cười trừ: ``Khoảnh khắc mà mình chạm mặt với những hồn ma ở đó\dots\ cảm giác vẫn còn quá thật, thậm chí là đến tận bây giờ. Mình chưa dám tin, nhưng mỗi lần nghĩ nhớ lại khoảnh khắc hôm đó là lại thấy nổi da gà. Đặc biệt là cô bé có đôi mắt xanh ấy.''

``Cô bé mắt xanh ư?'' tôi vừa kịp mở miệng.

\textbf{*BẠP!*}

Shion-san bên cạnh đập bàn, đứng phắt dậy, tay chống hông, ánh mắt sục sôi: ``Chờ đã! Đó chính là dấu hiệu một thực thể bị phong ấn đang thức tỉnh! Ta dám chắc trong lòng rừng có khối ma lực nguyên thủy. Ngươi thấy đấy, các dòng chảy linh hồn đều hướng về điểm ấy--''

``Cậu định triệu hồi nó lên rồi ký khế ước đấy hả?'' Hotaru-san ở bên cạnh lên tiếng, giọng hơi châm chọc.

Cô nàng tóc trắng-bạc hất tóc, không hề giận: ``Tất nhiên! Cơ hội ngàn năm một thuở để ta chứng minh thực lực, sao ta có thể bỏ qua!''

Azuki-san lắc đầu, quay sang tôi, giọng nhỏ nhưng đủ cả lớp nghe rõ: ``Mình không bàn chuyện triệu hồi. Mình chỉ thấy lần trước bọn mình mới chỉ lướt qua bề mặt. Cậu có muốn quay lại, lần này đi sâu hơn không?''

Tiếng xì xào bỗng im bặt. Những ai từng bước vào rừng đêm hôm ấy đồng loạt im lặng. Không còn là trò đùa, cũng chẳng còn là trò thám hiểm cho vui. Ánh mắt họ nhìn nhau, như thể cùng nghe thấy một bản nhạc chỉ riêng khu rừng mới phát ra.

Suzuki khẽ kéo tay áo tôi, giọng thì thào: ``Onii-sama\dots\ họ cũng đang cảm nhận được tiếng gọi ấy, đúng không?''

``Ừ, có vẻ là vậy rồi,'' Kaigou-chan cất tiếng, mắt vẫn dán vào bản đồ nhỏ đang trải trên bàn, ``không chỉ nghe, mà còn đang bị \emph{kéo vào}.''

Tôi gật nhẹ. Không chỉ riêng hai chúng tôi; mọi người từng đặt chân vào khu rừng Aogami đều đang bị kéo về phía ấy bởi một thứ không tên.

``Được rồi, đồng minh!'' Shion-san vỗ bàn, mắt lóe như vừa mở khóa ấn chương thất truyền. ``Chiến dịch \textbf{Chiến dịch Hồi Cổ Lần II} chính thức khởi động!''

``\dots Tên nghe nguy hiểm quá đấy,'' tôi buột miệng.

Cô nàng phù thủy tự phong nâng cằm, đôi mắt tím biếc chìm trong bóng tối huyền bí. ``Kuon, đó là bởi lớp vỏ ngôn từ của kẻ phàm tục không thể chứa nổi hương vị \textit{đại hỗn loạn} sắp tràn về!'' Giọng cô bỗng trầm xuống, như thể gió âm phủ rít qua khe cửa. ``Lần trước ta chỉ đạp cổng, lần này, ta sẽ đối mặt với \textit{con mắt nguyên thủy} đang chớp giữa tán lá!''

``Tôi thấy cậu chỉ đổi mỗi cách gọi thôi,'' Sakura-san phản bác, ``có khác gì nhau đâu.''

``Không! Đó là sự khác biệt giữa \textit{lướt nhẹ trên mặt ảo} và \textit{đâm thẳng vào trái tim bóng tối}!'' cô siết tay, khí thế như kiếm sắp rút vỏ. ``Tinh thần chiến binh linh thiêng, thứ duy nhất có thể bảo vệ chúng ta khỏi bị \textit{quá khứ nuốt chửng}!''

Tôi thở dài ngán ngẩm. Đúng là hết thuốc chữa với cô nàng bị hoang tưởng nặng này.

Shion-san còn đang hăng máu phô diễn kế hoạch ``đại hỗn loạn'' thì một giọng cao vút cắt ngang, như ai gảy thúng úp vào đầu cả nhóm. ``Hể\~{}? Có trò gì vui vậy?''

Suzune-san. Cô ấy lại xuất hiện kiểu ``thêm vào sau khi in'': không tiếng bước chân, không hơi thở, chỉ thấy thân hình nhỏ nhắn đứng đó, mắt tròn xoe như đèn pin pin năng lượng mặt trời. Cô nhìn lần lượt từng người, chờ đợi một lời mời gọi.

``Có gì vui không? Cho mình tham gia với!''

``Bọn mình đang tính vào rừng,'' Azuki-san đáp ngay, không giấu nổi sự hào hứng. ``Phiêu lưu ban đêm luôn.''

``Hể\~{} Thật không!?'' cô nàng tóc vàng nhún gối, hai tay vung vẩy, cả người như con sóc vừa phát hiện thêm một kho hạt dẻ. ``Mình muốn! Cho mình đi với!''

Shino-san khẽ kéo tay cô bạn, giọng hạ thấp nhưng đầy cảnh báo. ``Cậu nên suy nghĩ kỹ một chút đi, Suzune-san. Nó sẽ không thú vị như cậu nghĩ đâu.''

``Eh\~{} nhưng mà--''

``Nghe lời mình đi,'' cô ấy nói nhẹ, nhưng ánh mắt xanh dương đó lại không hề đùa. ``Có những thứ mà mình nghĩ\dots\ cậu không nên thấy thì hơn. Sẽ có những thứ đáng sợ khiến cậu ám ảnh đấy.''

Suzune-san chớp mắt, nụ cười hơi cứng lại. Chỉ một giây, ánh sáng trong mắt cô như bị mây đen qua. Nhưng rồi cô lại bật cười, nhún vai: ``Mình biết mà! Nhưng nếu không thử, sao biết mình có hối hận hay không?''

Suzuki thở dài, định nói thêm thì tôi không kìm được mà xen vào mỉa mai: ``Cậu cứ thế là `cả thèm chóng chán' đấy.'

Câu nói vừa dứt, cô nàng ``Linh vật'' như bị ai bấm nút tạm dừng. Khóe môi cô chu lên, má hơi phồng: ``Kuon-kun xấu tính quá! Cậu nói cứ như mình chưa từng kiên nhẫn bao giờ ấy!''

``Không phải sao?''

``\textbf{Không phải!}'' cô phản kháng, nhưng đôi mắt xanh lá lại lảng đi, như tìm chỗ trốn cho lý lẽ của chính mình. Tôi biết mình đã trúng đích: cô ấy chính là kiểu người hễ nghe ``cấm'' lại càng muốn xông vào, chỉ để rồi sau đó ôm hận và kể lể ``thật ra cũng chẳng vui như tưởng''.

``Khu rừng\dots'' Một tiếng thì thầm nhỏ như lá rơi xen vào giữa bầu không khí đang nóng lên. Mari-san đứng lặng ở cuối lớp, cuốn sổ vẽ ôm chặt trước ngực, đôi mắt hai màu mở to như vừa nhìn thấy bóng ma quen thuộc. Cô ấy không nhớ mình đã từng bước vào khu rừng ấy - ít nhất là trong đời thực này - nhưng cơn sóng thân thuộc ập tới mạnh đến mức đầu gối cô run nhẹ. Mùi đất, tiếng côn trùng rít, cảm giác có ai đang nhìn xuyên qua tán thông\dots\ tất cả tràn về như thể vừa được gỡ khỏi một lớp bụi thời gian dày đặc.

``Tôi\dots\ tôi cũng muốn đi,'' cô nàng cất giọng, khản hơn bình thường. Cả nhóm chúng tôi quay lại. Cô nuốt khan, nhưng không rút lời. ``Tôi không biết tại sao, nhưng\dots\ tôi cảm thấy nơi đó đang gọi mình.'' Cô siết chặt cuốn sổ, những trang giấy vẫn còn ẩm hơi chì mờ đục. Trong đầu cô, bức phác thảo khu rừng cô vẽ vẫn giật giật như muốn thoát ra ngoài, như thể chỉ cần bước vào rừng thật, nét bút sẽ tự tìm đường về đúng vị trí đứt gãy của ký ức.

Shion đập tay xuống bàn, reo lên: ``Tuyệt! Thêm một đồng minh nữa!''

Suzuki khẽ nhíu mày, giọng không giấu được sự e ngại: ``Cậu chắc chứ? Mình thấy tâm lý của cậu hiện tại không ổn đâu.''

``Nếu cậu không đủ can đảm thì không cần phải đi đâu,'' tôi tiếp lời.

``T-Tôi ổn mà,'' Mari-san ngẩng lên, giọng run nhẹ nhưng không chùn bước. ``Tôi cần\dots\ câu trả lời. Nếu khu rừng đang gọi, thì đó là cách duy nhất để tôi biết tay mình vẽ ra điều gì, hay là điều gì đang vẽ lấy tôi.''

Tôi gật, nhìn cô nàng tóc nây ấy. Ánh mắt cô không còn hoang mang như lúc nãy; thay vào đó là một thứ kiên định mới chớm, như thể cô vừa nhận ra mình là mảnh ghép thiếu của bức tranh đêm nay. Khu rừng đang chờ, và cả hai chúng tôi đều biết: lần trở lại này, không ai được phép làm khán giả.

``Thôi, bỏ chuyện linh tinh qua một bên.'' Shion gõ nhẹ lên mặt bàn, tiếng *CỘP* khô khan đủ kéo mọi ánh nhìn về phía cô. Từ trong cặp da đen nhánh, cô rút ra một xấp giấy đã ngả vàng, kèm vài tấm ảnh chụp lén hạt đậu nhỏ, trông như bằng chứng của một vụ án chưa bao giờ được công bố.

``Ta đã thu thập manh mối từ trước khi bình minh kịp thức dậy.''

Sakura-san nghiêng đầu, lông mày rướn cao: ``Từ lúc nào vậy, thưa `đại thám tử' Shion?''

Shion-san khẽ vuốt mái tóc trắng-bạc, đôi mắt lóe lên tia sáng của kẻ đang nắm giữ bí mật vũ trụ. ``Đương nhiên là sau khi ta xé toang tấm màn ngăn cách thế giới này với Thư Viện Tối Thượng, nơi những trang nguyền rủa được giam cầm bởi bảy tầng phong ấn thời khai thiên--''

``\dots Nói tiếng người giùm đi.'' Kaigou-chan hừ một tiếng, rõ rằng khó chịu.

``\dots'' cô nàng Phù thủy xoay nhẹ cổ tay, cánh tay phải như vừa bị trúng phải lời nguyền im lặng. ``\dots Ở một kệ tủ của CLB huyền bí. Có một ngăn kéo bị kẹt, ta dùng thìa gõ cho nó nhả ra.''

Cô chị lớn thở dài, tiếng gió lạnh lùng như muốn đóng băng cả căn phòng. `Đấy. Nói thẳng từ đầu có phải tốt hơn không. Đao to búa lớn làm gì.''

Shion-san quay mặt đi, giả vờ ho khan như thể vừa bị bụi thời không làm sặc cổ. Cô lật tài liệu, giấy ken két, như thể mỗi trang đều rên rỉ vì bị lôi ra khỏi quá khứ. ``Thôi, bỏ qua hành trình. Giờ hãy nhìn vào bản đồ dòng linh hồn - thứ sẽ dẫn chúng ta vào tim của `con mắt' đang chớp giữa rừng sâu.''

Tôi khom lưng, đưa mắt xuống tập hồ sơ đã ngả vàng. Mực xưa loang lổ, có dòng bị sau này ai đó ghi thêm bằng loại bút khác, rõ ràng đây là thứ được truyền tay, vá víu qua bao thế hệ. Trang đầu tiên chỉ còn dòng tiêu đề mờ nhạt:

\txtbx{
  {\vnmsans "Ghi chép dân gian - khu rừng phía bắc Aogami (cuối thế kỷ XIX - đầu thế kỷ XX)"}}

``\dots Lâu đến tận hồi đó cơ à,'' tôi lẩm nhẩm.

Shion-san khẽ gật, mái tóc trắng-bạc rung nhẹ như bị gió vô hình thổi. ``Lời nguyền ấy đã ngủ yên hơi lâu rồi.''

Cô lật tiếp. Những hàng chữ hiện ra, nhỏ nhưng sắc:

\txtbx{
  {\arial "Người dân tin rằng các vị thần trên núi cần được 'xoa dịu' định kỳ. Vào năm mất mùa hay dịch bệnh, họ tổ chức nghi thức hiến tế."}}

Dòng tiếp theo khiến không khí như nghẹn lại: ``Đối tượng là trẻ em.''

Khỏi cần nhìn cũng biết Suzune-san đang rụt cổ. Kaigou-chan im lặng, hai ngón tay siết nhẹ viền áo. Chỉ còn tiếng giấy sột soạt, như có bàn tay vô hình lật từng trang.

\txtbx{
  {\arial "Nghi thức diễn ra trong một khu vực sâu trong rừng, ở nơi được gọi là `vùng đất cấm.' Tại đây, đứa trẻ sẽ được đặt vào hố đất được đào sẵn, chôn sống cùng bùa ngải, bình phong gỗ và hạt lúa rang."}}

Tôi thấy gáy mình lạnh ngắt. Trên giấy, mực in mờ nhưng vẫn ngửi thấy mùi đất ẩm.

\txtbx{
  {\arial "Người ta tin rằng linh hồn của đứa trẻ sẽ trở thành 'sứ giả', mang theo lời cầu nguyện của con người đến các vị thần."}}

Shino-san khựng lại, đôi mắt xanh dương như bị thôi miên. ``Nhưng,'' cô thì thầm, ``có ghi chép khác nói\dots\ linh hồn ấy không bao giờ rời đi.''

Tôi nắm chặt tay, móng chọc vào lòng bàn.

\txtbx{
  {\arial "Chúng vẫn ở lại. Trong bóng tối, dưới lòng đất, giữa rừng già. Oán niệm tích tụ theo năm tháng. Không được giải thoát. Không được lắng nghe."}}

Dòng cuối cùng của đoạn được viết bằng mực mới hơn, nét bút run run:

\txtbx{
  {\arial "Những đứa trẻ không rời đi. Chúng chờ đợi một điều gì đó."}}

Im lặng bao trùm. Có tiếng lá ngoài hành lang va vào cửa kính, như bàn tay nhỏ gõ từ bên kia thời gian.

Shion-san hít một hơi dài như thể đang hút cả linh hồn của bóng tối vào lồng ngực, rồi rút ra một tờ giấy cuối cùng. Trên đó là bản phác thảo khu rừng bằng những nét chì mờ như bị thời gian cắn gặm, và một vòng tròn đỏ au, đỏ đến mức như máu vừa mới khô, được khoanh lại bằng mực đậm đặc hơn mọi chỗ.

\txtbx{
  {\arial "Vùng đất cấm - tránh tiếp cận sau khi mặt trời lặn."}
}

Tôi nhìn chằm chằm vào vòng tròn. Nó nằm sâu trong rừng, cách ngôi đền Aogami chừng bốn trăm mét về phía tây bắc, nơi mà Saya-san và Touka-san đã chỉ tay trong đêm vào rừng lần trước.

Kaigou-chan thì thào:  ``\dots Chính xác là chỗ đó.''

Sakura-san im lặng, hai ngón tay siết chặt viền áo, mắt cô lóe lên tia sợ hãi hiếm khi thấy.

Tôi không nói gì, chỉ cảm thấy gáy mình lạnh toát, như thể có ai đó vừa thổi một hơi từ phía sau.

Azuki-san cúi xuống, nhíu mày: ``Còn hiện tượng thì sao?''

Shion nở nụ cười của kẻ vừa mở được cánh cổng địa ngục. Cô lật tiếp tờ giấy, giọng trầm xuống như thể đang đọc một bản thánh ca bị nguyền:

\txtbx{{\arial "Nghe tiếng trẻ con hát vào ban đêm. Thấy bóng nhỏ di chuyển giữa các thân cây. Đôi khi, có người cảm thấy như bị ai đó nắm tay kéo đi."}}

Cô nàng Phiêu lưu giật mình: ``\dots Tiếng trẻ con hát?`` Cô lẩm bẩm, mắt mở to: ``Honoka-san cũng từng nói vậy, đúng y chang.``

Shino-san khẽ lùi lại một bước, giọng run run: ``Mình\dots\ mình không thích đoạn cuối chút nào.``

Sakura-san gật nhẹ, tay đã lục tìm lá bùa trong cặp: ``Cảm giác bị kéo tay\dots\ không phải là ảo giác đâu. Có thứ đó thật.``

Azuki-san cười trừ, nhưng nụ cười hơi cứng:  ``Tớ thì lại thấy khá thú vị đấy.``

Shion-san đập bàn, mắt lóe tia điên cuồng của kẻ vừa tìm ra chìa khóa vũ trụ: ``Ta cũng thế! Đó là tiếng gọi của `bên kia'! Các linh hồn nhỏ đang mời chúng ta vào vòng xoáy thời không!``

Suzune nghiêng đầu, ngón tay quay quay lọn tóc vàng: ``Ừm\dots\ mình thì\dots`` Cô ngập ngừng, rồi cười tươi như con sóc vừa tìm được hạt dẻ phát sáng: ``Nghe hơi đáng sợ, nhưng mà\dots\ nếu có thật thì chắc sẽ rất đặc biệt nhỉ?``

Đúng kiểu cô ấy. Sợ thì có sợ. Nhưng chỉ cần thấy ``thú vị`` là kiểu gì lao vào như con thiêu thân, mặc cho cái ``thú vị'' này không được vui vẻ cho lắm.

Kaigou-chan hạ giọng, mắt vẫn dán vào bản đồ cũ: ``Còn manh mối nào nữa không?''

Cô nàng Phù thủy tự xưng không đáp lại ngay. Cô khép mí mắt lại, như thể đang lắng nghe tiếng thì thầm từ một chiều không gian khác. Một luồng gió lạ chạy dọc sống lưng tôi khi cô chậm rãi rút từ trong áo choàng đen một tấm ảnh. ``Tấm ảnh này,'' cô ngẩng lên, khóe môi nhếch một nụ cười bí hiểm, ``là do Uzuki bắt được khi ả đi săn `bướm đêm,' những tinh linh lỡ lạc vào cõi trần.''

Tôi nghiêng người. Bức ảnh mờ ảo: hoàng hôn tím tái, cành thông nghiêng ngả như bộ xương khổng lồ. Và ở trung tâm, có một bóng hình nhỏ nhắn, tóc tết dài quá gối, hai mắt phát ra thứ ánh sáng xanh lục không thuộc về thế giới này.

Suzuki bất giác nắm chặt tay tôi. Cô thì thào, giọng run như sợ gió nghe được: ``\dots Là cô bé đó.''

Không ai cần hỏi thêm. Trong căn phòng ngột ngạt, ba tiếng ấy đủ khiến tất cả hiểu: kẻ đứng đầu đám linh hồn trẻ con - một thiếu nữ chưa bao giờ lớn, kẻ đã đánh mất tên từ thuở khai thiên - vừa lộ diện.

Sakura-san vẫn lặng như tượng từ nãy, bỗng khựng lại. Đôi mắt cô dán chặt vào bức ảnh, dường như nhận ra điều gì đó từ thời ký ức chưa từng mở lời.

Sakura nhìn chằm chằm vào tấm ảnh, môi mấp máy như thể vừa ghép được mảnh cuối của một câu đố cũ: ``Khớp rồi\dots\ đúng những điều tôi từng nói ngày trước.'' Tiếng cô nhỏ đến nỗi chỉ có bản thân cô mới nghe thấy, nhưng tất cả chúng tôi đều cảm được cái run nhẹ trong lời ấy.

Phía sau lưng tôi, Mari-san siết cuốn sổ vẽ vào ngực, đầu ngón tay trắng bệch. ``Khu rừng\dots\ tôi chưa từng bước vào, nhưng sao mọi thứ lại quen thuộc đến thế.'' Câu nói vừa dứt, không khí quanh cô như nổi sóng, một làn gió vô hình thổi tung những trang giấy trong lòng bàn tay cô, dù cửa sổ vẫn đóng kín.

Shino-san im lặng, mắt hướng về nơi xa xăm. Tôi bắt gặp đôi vai cô khẽ run, một kiểu run của người vừa gặp lại một giấc mơ đã cố quên. Một thứ quen thuộc không được phép tồn tại trong hiện tại.

Tôi đẩy ghế đứng dậy, tiếng gỗ kêu ken két như thúc dứt sự im lặng.

``Dù sao thì,'' mọi ánh mắt quay về phía tôi, như thể tôi vừa kéo màn bước ra sân khấu, ``tôi định vào rừng kiểm tra ngay đêm nay.'' Không ai hỏi lý do, bởi họ đều biết: câu trả lời đang nằm sâu trong bóng tối phía tây bắc ngôi đền, và nếu chúng tôi không đi, chính bóng tối sẽ bò ra ngoài tìm chúng tôi.

Kaigou-chan gật đầu, mắt xanh lấp lánh cái nhìn của kẻ đã cân nhắc xong mọi rủi ro. ``Tôi đi cùng.''

Bên cạnh, Suzuki khẽ nắm tay áo tôi, ngón tay em lạnh toát. ``\dots Em cũng vậy.'' Không cần nhiều lời, tôi hiểu: dù có mọi thứ xảy ra như thế nào, chúng tôi cũng sẽ chịu khổ cùng nhau.

Không cần thêm lời dài dòng, cũng chẳng cần ai giơ tay biểu quyết. Chúng tôi đã hiểu: kẻ muốn đi sẽ đi, kẻ còn do dự cũng sẽ bị kéo theo. Thế là, nhóm được sinh ra trong im lặng, một cấu hình rời rạc, không ai giống ai, cũng chẳng ai bảo ai.

Tôi - người vừa tuyên bố cuộc hành trình.
Kaigou-chan - tay thám hiểm lặng lẽ, mắt xanh như đang đo gió.

Suzuki - em gái bé nhỏ, ngón tay vẫn run, nhưng nắm áo tôi không buông.

Shino-san - cô gái từng bị mắc kẹt giữa những vòng lặp miệt mài gõ nhịp lên cuốn sổ, như đếm nhịp tim rừng.

Sakura-san - cô tiên núi giấu bùa trong túi, mắt lúc nào cũng như đang cân electron.

Shion-san - tự xưng thủ lĩnh, giọng vang như loa phóng thanh, tay cầm tập hồ sơ ngả vàng.

Azuki-san - cô nàng phiêu lưu và không hề sợ hãi.

Mari-san - người cầm sổ vẽ, mắt hai màu mờ đi, như người vừa tỉnh giấc ở một đời khác.

Và Suzune-san - thêm vào sau khi in, mắt tròn vo, chỉ cần nghe ``có trò'' là nhảy vào.

Một bầy quái thú lạc đàn, chưa kịp ráp lại đã muốn xông vào đêm.

Azuki gõ bàn, cười toe toét: ``Vậy mình đi lúc nào? Mai? Mốt? Cuối tuần?''

Tôi không do dự mà đáp ngay: ``Tối nay. Khi trời bắt đầu chập choạng tối, ngay sau giờ học.''

Suzune giật mình, lông mi vàng khẽ rung: ``Nhanh vậy!? Để mình còn về nhà chuẩn bị đồ--''

``Không cần đèn. Có hết đồ nghề rồi. Với lại\dots'' tôi đáp, giọng thấp hơn mức cần thiết, ``đêm là lúc chúng ta, và cả chúng nó đều nhìn rõ nhất.''

Câu cuối tôi nuốt vội, nhưng Kaigou-chan đã nghe thấy. Cô gật nhẹ, mắt xanh lóe lên tia ``đã biết''. Suzuki siết tay tôi thêm một nấc, không nói, chỉ thở.

Không ai hỏi thêm. Tiếng chuông vào lớp reo, như mũi tên cắt dây neo. Chúng tôi tan ra, mỗi người mang một phần bóng tối trong lòng, chờ đến giờ điểm.

Tôi quay ra cửa sổ. Dãy thông xa xa vẫn đứng im dưới nắng chiều tà, nhưng tôi biết khi mặt trời tắt, chúng sẽ thay áo. Không còn là cây cối bình thường, khu rừng giờ là một con mắt khổng lồ khép hờ, tròng đen ẩn sau hàng mi lá kim. Nó đang chờ những kẻ đủ can đảm bước vào, để viết tiếp những trang câu chuyện còn dở dang từ trăm năm trước.

\ct

Kế hoạch vừa được lập nên, cả bọn liền giải tán nhanh chóng như thể chẳng hề có cuộc bàn bạc nào xảy ra. Không ai dám nán lại quá lâu, cũng chẳng ai thốt ra hai tiếng ``khu rừng'' to tướng. Nhưng bầu không khí sau buổi họp chớp nhoáng đó lại nặng hơn bình thường, như thể mỗi người đều mang theo một chiếc đồng hồ riêng đang đếm ngược đến lúc mặt trời tắt.

Buổi chiều học thế là trôi trong màu lạ. Tôi vẫn cầm bút, vẫn ghi, nhưng những con chữ cứ nằm rạp trên trang giấy như thể sợ bị ai đó đọc thấu. Kaigou-chan đáp bài đến độ chính xác làm người khác nghẹt thở; cô ấy dường như muốn chứng minh rằng thế giới vẫn quay đúng trục, dù ai cũng biết trục ấy đang nghiêng. Suzuki ngồi ngay ngắn, mắt dán bảng, nhưng đuôi mắt lại cứ vụt ra cửa sổ mỗi khi có lá rơi, những chiếc lá như lá thư rừng gửi tới, chỉ có em ấy mới đọc được.

Mari-san không chạm bút. Cô mở cuốn sổ, lật đi lật lại, trang trắng nào cũng như tấm gương soi thấy bóng quá khứ. Đôi lúc cô dừng lại, hơi thở gấp, như thể nét vẽ sắp nhảy ra khỏi giấy và bám vào cổ tay cô. Shino-san gần như tàng hình: cô ngồi khép nép, tay siết cuốn sổ nhỏ, môi mấp máy đếm nhịp im lặng. Suzun-san thì vẫn tươi, vẫn cười, nhưng có những khúc quanh cười đột ngột bị cắt nghĩa bằng khoảnh khắc trống rỗng, như ai đó vừa bấm nút tua điện thoại, quên mất mình đang ở đoạn nào.

Giờ nghỉ trưa, Azuki-san và Shion-san lẩm bẩm bên góc lớp. Tiếng họ nói nhỏ như muỗi kêu, nhưng vẫn lọc được vài mảnh: ``đường vòng'', ``tiếp cận từ hướng đông'', ``đừng để bị kéo tay''. Những cụm từ ấy lẫn trong tiếng cười giòn giã và tiếng nói râm ran của mọi người, như mảnh vụn thủy tinh ẩn sâu trong bánh ngọt mà ai không cẩn thận sẽ nuốt phải vào lòng.

Tôi không làm gì cả. Tôi chỉ ngồi, để thời gian chảy qua kẽ tay, và chờ. Chờ bóng tối gõ cửa, chờ khu rừng thở dài thứ tiếng gọi đã níu chân mỗi đứa chúng tôi từ buổi sáng. Khi tiếng chuông cuối cùng vang lên, tôi biết, màn đêm sẽ không còn là bóng tối bình thường nữa, mà là một tấm màn ướt đẫm mực, sẵn sàng in dấu chân chúng tôi lên những trang giấy chưa từng được viết.

Tiếng chuông vừa dứt, lớp 1-C như bị xả van, học sinh ùa ra ngoài. Lớp học dần thưa bóng học sinh. Chúng tôi thì ở lại, mỗi đứa một lý do: người làm công tác trực nhật dù không có ca trong tuần, người thì giả vờ quên đồ để giết thời gian, người thì vào thư viện hoặc câu lạc bộ để sinh hoạt.

Bóng nắng từ từ lặn từ phía tây, nhuốm viền cửa sổ màu cam nhạt. Khi dãy nhà đổ bóng dài trên sân, chúng tôi mới lần lượt khép lại cửa lớp, bước ra hành lang, hướng về cổng sau, nơi không ai để ý.

Không cần gọi tên, chỉ cần ánh mắt. Tám bóng người lặng lẽ tụ lại dưới gốc cây bàng, bóng cây khổng lồ trải dài trên nền đất như một tấm màn đen nhung thấm đẫm mùi thoang thoảng của cỏ cây. Gió chiều tà lùa qua kẽ lá, rì rào như tiếng thì thầm của kẻ đã đợi sẵn từ lâu. Mặt trời chỉ còn là một vệt cam nhạt ở rìa tây, đủ để phơi bày những đường viền gai góc của hàng thông phía xa, nhưng không đủ để xua tan cái lạnh đang bốc lên từ lòng đất.

``\dots Đi thôi, mọi người.''

Chúng tôi vừa bước khỏi bậc cổng cuối cùng, đôi giày còn chưa kịp in dấu lên lớp lá khô của rừng rậm thì một giọng nói mềm mại nhưng đủ sức kéo cả hàng người dừng lại đã cất lên từ phía sau.

``Các em định đi đâu thế?''

Yuko-sensei dựng thẳng thân hình mảnh khảnh dưới tán bàng, ánh đèn đường phủ lên mái tóc vàng sáng. Cô không chau mày, cũng chẳng cười, chỉ khoanh tay như thể đang đợi một câu trả lời vừa đủ thuyết phục cho hành động sắp tới của chúng tôi.

Tôi cố gắng để giọng mình trở thành làn gió, rồi cất giọng đáp bình thản nhất có thể: ``Thưa cô, bọn em chỉ đi dạo một vòng quanh khu đền, kiểm tra lại mấy tin đồn học sinh trong trường ạ.''

``Tin đồn?'' cô nghiêng đầu, lọn tóc lòa xòa che khuất một phần trán. ``Tin đồn nào cơ?''

Câu hỏi bật ra quá nhẹ, như thể cô đang đọc thành tiếng dòng chữ trong đầu mình chứ không hỏi chúng tôi. Tôi cảm thấy bàn tay đang cầm đèn pin sau lưng bỗng ướt đẫm mồ hôi.

``Dạ\dots\ chỉ là chuyện vớ vẩn thôi ạ. Cô cũng không cần quá để tâm đâu,'' tôi cố nhún vai thật tự nhiên, nhưng cơ vai cứng đờ như khúc gỗ.

Shion-san hé miệng, định bổ sung thêm chất liệu huyền bí cho câu chuyện thì hai bàn tay từ hai hướng nhanh như cắt bịt kín miệng cô. Kaigou-chan thì thầm đủ nghe: ``Cậu mà hé thêm một chữ, tôi bỏ cậu lại với cô đấy.'' Suzuki gật mạnh, ánh mắt như năm roi quất vào trán đối phương.

``Ưm\dots\ ưm!'' cô nàng Phù thủy giật người, mái tóc bạc phất loạn xạ trong gió.

Yuko-sensei nhìn bộ ba chúng tôi - một đứa cố tỏ ra vô tội, hai đứa còn đang khóa miệng bạn - rồi thở ra một hơi dài. Không phải kiểu thở của người vừa bị lừa, mà là kiểu thở của ai đó đã đọc trước đoạn kết, chỉ chờ xem diễn viên thoại sai chỗ nào.

``Kuon-san.'' Giọng cô như mặt hồ đột ngột tối sầm. ``Tôi đã từng nói với em rằng đừng nên để bản thân chịu đựng một mình mà nhỉ?''

Tôi cứng đờ, môi vừa mở đã khô khốc.

``Nếu có chuyện gì khó khăn\dots\ thì cứ đừng ngần ngại mà chia sẻ. Đừng để nó ăn mòn trong đầu.''

Tôi thấy mình đang chuẩn bị thành thật.

Nhưng, cô nhanh hơn. ``Tiếng trẻ con trong rừng, phải không?''

``\dots!'' Một luồng điện chạy dọc sống lưng. Cả nhóm như bị ai bấm nút tạm dừng.

Không chỉ bất ngờ; đó hoàn toàn là một câu đúng tới chí mạng.

Tôi nhớ lại buổi chiều hôm ấy, tại ngày thứ hai mà tôi theo học: bảng phấn còn vết xóa, ghế gập chỏng chơ, và cô đứng yên bên cửa sổ, nói khẽ về ``những âm thanh không nguồn gốc''.

Yuko-sensei thở ra, giọng trầm như đất ẩm: ``Tôi biết Aogami như biết lòng bàn tay: từng trận sạt lở, từng dấu chân chìm trong bùn, và cả những lời thì thầm về khu rừng không bao giờ ngủ.''

Kaigou và Suzuki liếc nhau, cùng nhớ cái ngày cô bước vào lớp 1-C, váy dài rung nhẹ, và đã cảnh báo tôi ngay từ buổi đầu. Cô chủ nhiệm chống hông, ánh mắt rõ ràng là không hài lòng với thái độ giấu giếm của tôi: ``Chẳng phải tôi đã từng nói như thế với em sao, Kuon-san?''

Không khí như tảng đá vừa lăn xuống, đè nghẹt lời từ chối trong cổ họng tôi.

``Vậy,'' cô chủ nhiệm tiến thêm một bước, đủ gần để tôi ngửi được mùi gió núi vẫn vương trên tóc cô, ``lý do thật sự khiến các em vào rừng là gì?''

Giọng cô không to, nhưng mỗi chữ đều in vào đáy mắt, không cho phép né tránh.  ``Đừng tưởng có thể lừa được tôi.''

Tôi cắn nhẹ môi. Ban đầu, tôi định giữ kín như bưng. Càng ít người biết, càng ít ai bị vạ lây.
Nhưng những ngày trôi qua, tiếng khóc trong rừng đã vượt qua bức tường phòng ký túc. Lời đồn lan nhanh hơn lửa cháy khô kể từ thời điểm đó. Học sinh bắt đầu nghe thấy, nhìn thấy, và\dots\ cảm nhận thấy có sự khác thường.

Việc giấu diếm thêm bây giờ đã không còn khả thi, và không sớm thì muộn thì mọi người trong trường cũng sẽ biết, theo cách này hay cách khác.

Tôi thở dài, hơi ấm phả vào không trời se lạnh. ``Đầu đuôi là thế này ạ\dots''

Tôi kể, không khai tất, nhưng đủ để cô hiểu: những linh hồn nhỏ đang mong chờ một điều gì đó trong suốt hàng trăm năm, những sự thay đổi dạo gần đây ở cả thị trấn này, và lý do chúng tôi phải quay lại đó.

Yuko-sensei không hề chen ngang hay cắt lời tôi nói, mà chỉ lặng lẽ đứng đó, để từng lời tôi rơi xuống như hạt mưa phùn trên mặt hồ. Không gian sau lời kết của tôi như bị ai bóp nghẹt, im lặng đến mức nghe được tiếng lá ngoài hành lang run rẩy. Tôi ngẩng lên, gặp đôi mắt xanh nước biển của cô. Ánh mắt không chút nào giật mình, không lẩn trốn, chỉ phẳng lặng như tờ giấy trắng đã đoán trước dòng chữ sẽ được viết lên.

``\dots Ra là thế.'' Cô nói rất nhỏ nhẹ, kèm theo một hơi thở dài như xả hết bụi trời đã đọng lại từ lâu. ``Tôi hiểu rồi.''

Một khoảng lặng nhỏ, đủ để cả bọn kịp nín thở. Rồi cô ngẩng cao, nhìn thẳng vào chúng tôi, không cho phép ai lảng tránh: ``Vậy tôi sẽ đi cùng.''

Azuki-san suýt nghẹn với quyết định có phần chóng vánh của cô chủ nhiệm: ``\dots Dạ?''

``Nếu có chuyện gì xảy ra, tôi sẽ chịu trách nhiệm.'' Nghe không giống thương lượng, chỉ là thông báo đã ký duyệt. Cảm giác của một sự dứt khoát như viên đá cuối cùng được đặt vào bức tường thành.

Suzuki siết tay áo tôi chặt hơn, ngón tay lạnh ngắt. ``Yuko-sensei\dots'' tiếng em nhỏ đến mức chỉ còn hơi ấm run run trên vai, không còn vẻ lo âu như lúc nãy, mà như cánh buồm vừa tìm được gió để trở về.

Shino-san khẽ cúi, cơ mặt giãn ra nhẹ nhõm: ``Cảm ơn cô ạ.'' Chỉ thế thôi, nhưng đủ để cả bọn hiểu: lần này, họ không đơn độc.

Tôi vẫn im. Không cần lời nữa. Trong lồng ngực, thứ căng thẳng vừa được tháo nút, tuôn ra thành dòng nước ấm chảy về cuối huyệt. Ít nhất, đêm nay, chúng tôi sẽ không phải tự mình cầm đuốc giữa bóng tối.

\ct

Những lời kể, những im lặng xen lẫn cái gật đầu của cô chủ nhiệm, tất cả hòa vào nhau cũng đủ để mặt trời kéo tuột xuống chân trời. Ánh sáng cuối cùng rút khỏi mái ngói, để lại bầu trời nặng trĩu màu chì.

Tôi khẽ lên tiếng: ``Đi thôi.''

Không còn ai phản đối, nhưng cũng chẳng ai lùi bước. Chúng tôi chỉnh lại dây cặp, bước qua cổng sau trường, hòa vào dòng gió lạnh đầu đông. Con đường mòn phía dưới chân đền Aogami hiện ra trong bóng tối, mỏng như sợi chỉ đỏ nối giữa hai thế giới.

Bước chân ai cũng nhẹ, như sợ đất dưới gót giày vỡ tan. Không tiếng nói, không cười đùa, chỉ nghe thấy tiếng lá khô rơi rạo rực bên tai. Trong cái im lặng ấy, tôi bỗng thấy lòng bàn tay mình ấm lạ - có lẽ Suzuki đã nắm chặt từ lúc nào. Cảm giác quen thuộc ấy lại trỗi dậy: rằng phía trước, nơi tán cây đen như bức tường thày, có thứ gì đang giang tay đợi chúng tôi từ rất lâu rồi.

Rừng Aogami vào buổi tối khác hẳn so với ban ngày.

Ngay khi chúng tôi bước qua ranh giới giữa con đường đất và tán cây đầu tiên, ánh sáng như bị nuốt chửng. Ánh sáng yếu dần, bị lọc qua từng lớp lá, cho đến khi chỉ còn lại một thứ gì đó mờ đục, lưng chừng giữa ngày và đêm. Gió thổi lạnh ngắt, mang theo mùi lá mục và hơi ẩm của đất vừa được xới lên; tiếng côn trùng cũng nhỏ dần, nhường chỗ cho những âm thanh không thuộc về thế giới này: tiếng thì thầm như có ai đó đi ngay sau lưng, tiếng lá khô rung lên dù không một luồng gió, và thứ tiếng rít rất nhẹ, như xương thủy tinh va vào nhau.

Trên cao, tán thông xoắn lại thành mái vòm đen kịt, chỉ chừa những khe hở hình lưỡi liềm để ánh trăng non lọt xuống, bạc nhợt và run rẩy như vừa mới học cách chiếu sáng. Mặt đất bên dưới phủ một lớp lá khô dày, mềm và âm thầm như tấm thảm được dệt bằng tiếng bước chân của chính mình; mỗi bước đi, chúng tôi lún sâu đến mắt cá, để lại vệt hằn nhỏ xíu nhưng đủ cho khu rừng nhận ra: kẻ lạ đã vào lãnh địa của nó. Ở phía trước, sương muối bắt đầu bay lên, quấn quanh thân cây, làm nổi rõ những vệt sáng mờ ảo. Có lẽ chỉ là ảo giác, nhưng trông như bóng những đứa trẻ đang trốn sau thân cổ thụ, thở nhẹ lên gáy chúng tôi.

Tôi dẫn đầu đoàn, bước chân vô thức như có dây kéo. Không hẳn muốn thế, chỉ là dòng chảy đẩy tôi về trước; nếu không đi trước, ai sẽ gánh bóng tối thay họ?

Mắt tôi lướt nhanh, thu vào ký ức từng nhánh cây cong queo, từng mảng sương lơ lửng. Điều khiến tôi giật mình không phải âm thanh nào cụ thể, mà là khoảng lặng đang bò vào tai, như ai đó bịt trống nhĩ bằng bông gòn.

Kaigou-chan theo sát, bước đều đặn như đếm nhịp kim đồng hồ. Cô không lên tiếng, chỉ nhìn: gốc cây có vết xước bí hiểm, tán lá có khe hở hình mắt, không khí giữa hai thân cây rung lên như sợi tơ vô hình đang căng.

Suzuki khẽ kéo tay áo tôi, ngón tay lạnh như hạt mưa đêm. ``Vào rừng lúc này\dots\ em thấy lạnh cả trong xương.''

Tôi nắm tay Suzuki thật khẽ, đủ để cả hai cảm nhận được hơi ấm của nhau: ``Đừng sợ, có anh ở đây rồi.''

Phía sau, Shino-san và Sakura-san bước song song, nhịp bước đều như đang đếm thời gian. Cô nàng tóc nâu dừng lại, ánh đèn pin lướt qua bản đồ cũ: ``\dots Chỗ này không còn giống như lần trước.''

Nàng Vu nữ chạm vào thân cây, mũi giày vô tình xới lớp đất xốp. Nhìn xuống dưới chân, cô vô thức cất lời: ``Độ ẩm tăng nhanh, không khí nặng như bị ai đổ thêm chì.'' Cô khẽ day lòng bàn tay lên vỏ cây, nhíu mày: ``Có thứ gì đang tích tụ\dots\ đang nghe ngóng.''

Cách đó mươi bước, Azuki-san vẫn tươi cười, nhưng nụ cười không còn giòn tan: ``Cảm giác\dots\ như bước vào ruột một con quái vật đang ngáy.''

Shion-san phất tay, áo choàng đen rung phần phật: ``Chính xác! Chúng ta vừa vượt qua `tầng giao thoa thứ nhất'. Từ đây, mọi thước đều bị `thiên địa chúng' bóp méo. Không gian không còn là không gian, mà giờ nơi này đã trở thành--''

Kaigou-chan quay lại, cắt ngang: ``Nói ngắn gọn thôi.''

Cô nàng Phù thủy ho khan, đổi giọng trầm trọng: ``Ở đây\dots\ có vẻ bất thường.'' Cô còn lẩm bẩm đằng sau, như bất mãn: ``\hl{Làm người ta mất hứng\dots}''

Mari-san đi giữa hàng, ôm cuốn sổ vẽ sát ngực. Cô không nhìn đường, không nhìn trời, chỉ nhìn thẳng vào khoảng tối phía trước, như thể có một sợi dây vô hình đang kéo cô.

Suzune-san bám sát tôi. Ban đầu còn tíu tít như một đứa trẻ háo hức tới một vùng đất mới, nhưng giờ đây đã bắt đầu run lập cập: ``Lạnh\dots\ lạnh trong xương luôn ấy.''

Azuki-san bình tĩnh đáp lại: ``Nhiệt độ bình thường mà.''

Cô nàng ``Linh vật'' lắc đầu, giọng run: ``Mình cảm tưởng\dots\ có ai đang thổi vào hộp sọ mình vậy.'' Mồ hôi ứa ra trên trán.

Cuối đoàn, Yuko-sensei bước đều, không nhanh không chậm.
Ánh mắt cô lướt từng học sinh, từng tán cây, từng vệt sương, như đang so sánh thứ đang thấy với một bức tranh đã in sẵn trong đầu. Sự điềm nhiên của một người trưởng thành không hề bị lung lay trước màn đêm đang chảy xiết quanh chân.

Chúng tôi cắm đầu vào rừng, mỗi bước như đạp vào một lớp băng mỏng vừa mới đóng. Không khí dày lên, từng bước chân của chúng tôi càng thận trọng hơn, tiếng nói cười cũng bớt râm ran hơn. Từ ký ức lần trước, nếu tôi nhớ không nhầm, cái gốc cây cong queo kia, chỗ này đáng lẽ phải trống trải, hoặc có lối rẽ kia, trước đây chỉ là bụi rậm.

Không ổn. Cảnh vật như bị xáo bài, nhưng bộ bài vẫn còn thiếu vài lá.

Tôi vừa nghĩ vậy thì một bóng trắng lướt qua khe cây - chỉ kịp thấy góc áo - rồi tan vào màn sương. Không ai khác ngoài tôi trông thấy. Hay chính tôi cũng không chắc.

Mari-san dừng lại, thở gấp như người vừa chạy marathon trong mơ: ``Mình từng đứng đây\dots\ nhưng lúc đó trời nắng, và không có mùi lá ủ.'' Cô ôm cuốn sổ sát ngực, như thể trang giấy sắp nhảy ra ngoài.

Sakura-san chậm bước, mũi chân vẽ một vòng cung trên đất mềm. ``Linh khí không tụ thành dòng mà rời rạc thành nhiều mảnh, nhiều chủ thể, nhưng dính nhau bằng một mạng nhện không nhìn thấy.'' Cô khẽ cúi, hai ngón tay búng ra một tia sáng xanh lục rồi tắt ngúm. ``Có gì đó mờ ám ở đây.''

Yuko-sensei đứng cạnh, gió thổi tóc cô bay ngược như ngọn cờ lính. ``Giống như một hệ sinh thái bị lệch cân bằng\dots'' cô thì thầm, `` nhưng thay vì năng lượng sinh học, thì là--''

``\dots linh hồn,'' cô nàng Vu nữ hoàn thành vế câu. Cả hai cô trò nhìn nhau, ánh mắt của họ cũng đủ để hiểu chuyện gì đang xảy ra.

Tôi giơ cổ tay. Mặt đồng hồ hiện ba chấm nhỏ chớp đỏ, rồi dòng chữ trượt ngang: ``Không phải một.'' Dừng lại nửa giây, thêm dòng thứ hai hiện ra:  ``Rất nhiều. Đếm không xuể.'' Không một lời nói, không tiếng méo mó, chỉ hai câu in nhỏ cụt lủn, đủ cho chúng tôi đọc xong rồi biến mất, như thể Game cũng sợ inh tai kẻ lạ.

Suzune-san bỗng dừng phắt, mắt tròn xoe dán vào khoảng không trước mặt: ``Cái\dots\ cái gì thế này\dots?''

Cô lùi lại, gót chân chậm rãi lướt trên lớp lá mục, tiếng khô rạo rào như xương khẽ va. Một bước, rồi thêm một bước nữa, vai cô chạm phải thân cây như tìm chỗ núp.

``Mình\dots\ mình không ổn chút nào\dots'' giọng cô nhỏ đến mức gió thổi qua là tan.

Azuki-san đứng gần, nụ cười đã gỡ xuống từ lúc nào. ``Mình đã bảo đừng theo mà.'' Cô không trách, chỉ thở dài, mắt vẫn dán vào chỗ trũng phía trước, nơi bóng tối như được gấp lại thành một lớp mỏng, đen đặc.

Shino-san gạt nhẹ tay cô bạn tóc vàng, giọng đều đều như dây thừng thả xuống giếng: ``Bình tĩnh. Chúng ta gần tới rốn rồi.'' Cô khép cuốn bản đồ lại, chỉ còn khe hở đủ để một ngón tay chạm vào dấu X mực đỏ. ``Thêm một khúc nữa, rẽ trái. Cố lên nào, mọi người.''

Chúng tôi lần theo khe hở giữa hai hàng thông, cành khô gõ vào nhau nghe như chuông gió của đám tang. Lối đi hẹp đến nỗi vai phải cọ sát vỏ cây, mùi nhựa thông lẫn mùi đất xới khiến cổ họng tôi rát. Một bước, hai bước. Ánh sáng như bị ai vặn nút, tắt ngóm. Không chỉ ánh sáng, mà cả thời gian cũng như bị kéo giãn ra thành từng sợi dây mỏng, mắc vào những cành thông gầy guộc. Gió thổi qua khe cốt, rít lên như tiếng ai thở dài từ đáy cổ họng sâu thẳm. Phía trước, màn sương bắt đầu dày lên, không phải sương đêm bình thường, mà là thứ sương được ủ từ hơi thở của những kẻ đã ngã xuống nơi này trăm năm trước. Mỗi bước chân đều in lên một khoảnh ký ức không thuộc về mình: tiếng khóc trẻ con vọng lại từ lớp lá khô, mùi máu tanh thoang thoảng trong không khí lạnh, và cảm giác có ai đang đi bên cạnh dù vai chỉ chạm vào khoảng trống.

Một vòng trũng mở ra trước mặt, nhỏ hơn sân bóng rổ nhưng đủ để mọi tiếng thì thầm trong đầu ta ùa về một chỗ. Đất ở đây lõm xuống như lòng chảo, lớp lá khô mỏng hơn, lộ những vệt xới cũ kéo dài từ tim tới vành. Không cần ai nói, tất cả đều biết: đây là điểm không động của la bàn, nơi kim ngừng rung vì sợ. Trên thành trũng, những cây thông cong vẹo như cột sống của kẻ đang quỳ, tán lá dày đến mức ánh trăng không thể xuyên qua, để lại dưới đáy một hố đen đặc quánh, nơi âm thanh bị nuốt chửng và ánh mắt bị mù. Không chỉ là la bàn: cả trái tim cũng ngừng đập vài nhịp khi bước vào đây, như thể cơ thể tự nhận ra đây là ranh giới không được phép vượt qua nếu còn muốn giữ lấy hồn mình nguyên vẹn.

Shino-san bước xuống trước tiên, mũi giày in hằn lên đất mềm. ``Chính chỗ này,'' cô thì thầm, như thưa với ai đang nằm dưới lớp đất kia.

Sakura-san lập tức quỳ gối, trải tấm bản đồ cũ ra giữa lòng trũng. Ngón tay cô lướt nhanh, so từng đường gân lá với đường gân đất, rồi gật nhẹ: ``Đúng là như vậy. Nghi thức đã diễn ra ở đây. Cột cờ không còn, nhưng dấu chân tròn vẫn in hình.''

Yuko-sensei bước xuống cuối cùng. Cô cúi xuống, bàn tay lướt nhẹ trên mặt đất như vuốt tóc một đứa trẻ đã ngủ. ``Đất này đã bị xáo nhiều lần, không phải do mưa gió.'' Giọng cô trầm, đều, như đọc bản ghi chép của chính lòng đất. ``Là tay người. Họ đào, lấp, rồi lại đào, như thể cố xóa một dòng chữ nhưng mực đã thấm quá sâu.''

Cô đứng dậy, phủi tay, nhìn quanh vòng trũng. ``Nếu tài liệu cũ không nói sai, đây từng là nơi họ chôn những đứa trẻ\dots\ không phải một, mà là cả một bầy.''

Không ai đáp. Gió chợt ngừng. Lá cũng thôi rạo rực. Chỉ còn tiếng tim ta đập, như đang đếm từng linh hồn nằm dưới kia, và chưa chắc đã ngủ yên.

Gió bỗng nín thở, không phải lụi dần mà như bị ai bịt miệng. Lá ngừng xào xạc, côn trùng ngật ngưỡng, cả khu rừng trở thành hồ chì đặc. Tiếng nói cất lên từ mép hố, mỏng như lá khô, nhẹ như hơi sương, nhưng lạnh đến mức khiến đèn pin trong tay tôi rung lên.

``\dots Vậy là\dots\ mọi người đã quay lại đây thật.''

Tôi quay ngoắt lại. Giữa đám sương đen đặc quánh đó, một bé gái bất chợt lộ ra - lần thứ hai tôi thấy em, nhưng lần đầu với những người còn lại. Em chỉ đơn giản là hiện diện, như thể màn đêm vừa nhấc một lớp voan mỏng lên để lộ khuôn mặt thật. Vẫn mái tóc đen tết đuôi sam thấp thoáng dưới vành nón cổ, vẫn vạt kimono nâu sẫn thời Nara phủ kín mắt cá, viền áo thêu họa tiết hoa cúc mất màu, và vẫn đôi mắt xanh lục như hồ nước ngọt giữa rừng sâu, trong veo, bình thản, và lặng lẽ gắn thẳng vào chúng tôi trong đêm vào rừng đầu tiên ấy.

Azuki-san thốt lên trước tiên, giọng hạ thấp đến mức chỉ còn là luồng khí ấm: ``L-Là cô bé ngày hôm đó\dots''

Shion-san giật mình, tay siết chặt cổ áo choàng như thể sợ từng sợi vải bị gió đêm cắt đứt: ``Không có linh khí\dots\ mà vẫn đứng thẳng. Đây không phải u linh thông thường.''

Suzune-san lùi về phía sau một bước, mắt tròn xoe, tiếng thở gấp thành từng đợt: ``V-Vậy là có ma thật ư\dots?''

Yuko-sensei đứng ngoài rìa vòng trũng, môi mím lại thành đường thẳng. Cô không rút lui, cũng không tiến thêm. Cô chỉ khẽ nghiêng đầu, đôi mắt xanh nước biển đo lường từng centimet không khí quanh em bé như đang đọc một bản nhạc thiếu nốt.

Tôi hít đầy hơi lạnh, tiếng nói vang ra khô khốc, nhưng bình tĩnh như người lớn nói chuyện với trẻ con: ``Bọn anh muốn biết chuyện đã xảy ra với em, và\dots\ tại sao mọi thứ lại trở nên như vậy.''

Cô bé không trả lời. Đôi mắt lướt từng khuôn mặt, không oán giận, chỉ lặng lẽ đo lường, như thử xem ai đủ trọng lượng để đáng tin.

``Nơi này\dots\ không dành cho người sống,'' em ấy nói nhỏ nhẹ, giọng khẽ vang lên giữa khoảng trống mờ ảo của rừng đêm. ``Mọi người sẽ không hiểu được đâu.''

Shino-san bước lên, giọng đều như dây căng: ``Chị từng mắc kẹt ở vòng lặp. Chị hiểu cảm giác đó.'' Cô bé hồn ma ấy nhìn cô lâu hơn, như soi gương. Tôi cảm thấy có gì vừa nứt ra giữa hai thế giới.

``Tôi cũng vậy,'' Kaigou-chan nói khẽ.

``Và\dots\ chị nữa,'' Suzuki nắm tay tôi, giọng run nhưng rõ.

Từng người lên tiếng, không cần bàn trước, như dòng nước tự tìm khe đá. Em bé vẫn im, nhưng ánh mắt đã không còn đóng cửa.

Mari-san khẽ cau mày, giọng khàn như giấy ráp: ``Shino-san? Ý mọi người là sao vậy?'' Suzune-san co rúm lại, ngơ ngác nhìn chúng tôi nói mà không hiểu gi: ``Có chuyện gì xảy ra với mọi người sao?''

Shino-san lắc nhẹ đầu, đôi mắt vẫn dán vào khoảng tối trước mặt, chỉ đáp gọn: ``Không tiện giải thích đâu.''

Cô bé giữ im lặng, đôi mắt xanh lục như mặt hồ phủ sương. Gió hết thổi, lá ngừng rạo rực, cả khu rừng như nín thở chờ câu nói đầu tiên từ em. Rồi giọng em vang lên, mỏng đến mức chỉ còn là rung động trong không khí:

``Mọi người\dots\ có khi nào đã từng bị bỏ lại không?''

Câu hỏi có gì đó hơi bất ngờ và có sức nặng như hòn đá chìm. Không ai trong chúng tôi đáp lời ngay. Tôi cảm thấy tay Suzuki siết chặt tay mình, ngón tay lạnh ngắt. Một khoảnh khắc im lặng kéo dài đến mức nghe được tiếng sương đọng trên tán cây rơi xuống.

``Có.''

Suzuki lên tiếng, tiếng nói run nhẹ nhưng rõ ràng. Tôi thấy vai em rung lên, như thể vừa tháo được chiếc mũi kìm gắn từ lâu. Câu trả lời của em không chỉ cho cô bé, mà cho cả những bóng đen đang lẩn khuất sau thân cây.

Cô bé nhìn cô nàng mắt hai màu xanh-đỏ lâu hơn một chút, đôi mắt không còn là tấm gương đóng kín. Một khe hở nhỏ vừa mở ra giữa hai thế giới.

Shion-san bước lên, giọng đều đều như dây thừng thả xuống giếng: ``Để ta thử một thứ.'' Cô mở cuốn tài liệu cũ, lật nhanh như người tìm lại trang giấy đã rách trong ký ức. Dừng lại ở một trang có góc giấy gấp nếp, cô hô lớn như đọc thần chú:

\begin{center}
  ``Trong đêm không trăng, khi tiếng gió ngừng thổi,
  đứa trẻ được đưa vào lòng đất,
  như một lời cầu nguyện sống.''
\end{center}

Tưởng rằng đó chỉ là một cách mà cô nàng Phủ thuỷ tự xưng thử để ra oai, nhưng rồi\dots

Không gian rung lên, không phải rung động của gió, mà là rung của thời gian. Ánh sáng từ đèn pin bỗng méo mó, như bị ai bóp méo qua lăng kính cũ. Tiếng lá khô dưới chân không còn rạo rực, mà thành tiếng bước chân của người đi trước chúng tôi, những người đã từng đứng đây, cách đây trăm năm. Màn đêm như bị xé toạc, để lộ những hình ảnh mờ ảo: dáng người mặc áo trắng đang cúi xuống đất, tiếng khóc trẻ con vang lên từ lòng hố, và những bóng đen đứng thành vòng tròn, miệng lẩm nhẩm những câu thần chú đã bị lãng quên. Không khí nặng trĩu mùi máu tanh và hương trầm đốt ẩm ướt. Mặt đất dưới chân chúng tôi như mềm ra, như thể vừa được xới lên để chuẩn bị đón một đợt ``khách'' mới. Từng bước chân của chúng tôi giờ đây không chỉ in dấu trên lá khô, mà còn đè lên dấu chân cũ: những dấu chân nhỏ xíu, thấp thoáng in hình đôi guốc gỗ và đôi chân trần của trẻ con. Tiếng thì thầm không còn đến từ một nơi nữa; nó vang lên từ khắp mọi hướng, như một dàn hợp xướng vô hình đang hát bản điếu văn cho chính ký ức của khu rừng.

Tôi biết cảm giác này: không phải bị hút vào quá khứ, mà là quá khứ tràn ra, như nước vỡ bờ. Kaigou siết chặt tay cô, mắt nhìn về phía xa trong bầu không gian đang dần biến đổi. Suzuki nắm tay tôi chặt hơn, hơi ấm của em như ngọn lửa nhỏ giữa màn sương đen. Shino-san gập cuốn tài liệu lại, nhưng trang giấy vẫn rung nhẹ, như có ai đang đọc tiếp từ bên kia.

Ánh sáng lóe lên bao trùm cả vòng trũng, trắng xóa như tia chớp nổ giữa trời đông, xuyên thấu màn đêm và ký ức.

``\textbf{Mình\dots\ mình không nhìn thấy gì cả!}'' Suzune-san kêu lên, lùi dần về phía sau, mắt nheo lại.

``Không sao đâu, vẫn là ảo ảnh thôi,'' Azuki-san nói nhỏ, giọng bình thản như đang dạo trong công viên.

Sakura cắn môi, mắt dán vào tia sáng đang lan toả: ``Không phải ảo ảnh đâu. Những gì chúng ta sắp thấy\dots\ là thật đấy. Mọi người, cẩn thận nhé.''

Tôi siết chặt tay Suzuki và Kaigou-chan, giọng khàn đặc: ``\dots Ừm. Có lẽ\dots\ đây chính là điều mà em ấy muốn kể cho chúng ta.''

Lần này, ký ức không chỉ là hồi ức nữa. Nó đang kéo chúng tôi về phía trước, về một nơi đã từng là ``trước'' của ai đó.

\ct

Ánh sáng lúc nãy như tia chớp giữa trời đông giờ chỉ còn là lớp sương mỏng phủ xuống, dịu đến mức có thể nhìn thấy từng hạt bụi lơ lửng. Tôi bước một bước, mũi giày không còn lún vào nền đất ẩm, mà chạm phải nền đất cứng, được nghiền nát bởi bánh xe gỗ và đôi chân trần qua lại hàng trăm năm. Mùi khói bếp cay xè bốc lên từ những mái rạ đã ngả màu, hòa với hương gạo mới và mùi nhựa thông còn vương trên thớ gỗ chưa kịp khô. Trước mặt tôi, con đường đất rộng chừng hai mét dẫn thẳng vào trung tâm làng, nơi một cây mơ cổ thụ đang rụng hoa trắng lả tả như tuyết rơi.

Đám trẻ chạy tung tăng quanh gốc cây, đuổi bắt bướm bằng vòng hoa dệt từ cỏ lúa non. Chúng mặc áo nâu gai thô, tay áo buông dài, tùy tiện buộc bằng dải vải gai đã phai màu. Một cậu bé ngã sõng soài trước chân tôi, bật cười khanh khách rồi lồm cồm đứng dậy, rồi vui vẻ chạy tiếp, như thể chỗ tôi đứng chỉ là một khoảng trống trong không khí. Tiếng dép rơm lạo xạo, tiếng chuông gió treo dưới mái hiên leng keng đuổi theo gió thu nhẹ, xen lẫn tiếng gọi của một bà cụ từ hiên nhà gỗ, nơi tấm nứa đung đưa che khuất nửa khuôn mặt.

Tôi quay sang đồng đội. Sakura-san đã lùi lại phía sau một bước, mắt cô không rời khỏi mái hiên được chống bằng những thanh gỗ thông chẻ dọc, vạch đen kẻ trắng của muội đất thời Nara còn nguyên vẹn. ``Không có một dấu đinh sắt nào,'' cô thì thầm, ``toàn mộng âm dương, chèn khít như bài học đầu tiên của thợ mộc.'' Azuki-san vuốt nhẹ vào tấm nứa che cửa, ngón tay dừng lại trước hoa văn cúc - biểu tượng trường thọ và xua đuổi tà khí - rồi gật gù: ``Nơi này đúng là yên bình thật đấy\dots''

Shino-san đứng thẳng, mắt nhìn xa về phía đồi thông phía đông làng, nơi những dải vải trắng buộc vào cành - có lẽ là vật cúng tế thần linh - đang phấp phới trong gió. ``Chúng ta đang đứng trước khi lời nguyền được tháp vào xương đất,'' cô nói, giọng trầm đến mức dường như chỉ mình tôi nghe được.

Tôi không đáp. Tôi chỉ nhìn theo một bé gái đang ngồi thu lu trên bậc gỗ, hai tay ôm con búp bê làm bằng vỏ sò và rơm rạ. Em hát khẽ một điệu ru cổ, tiếng ``yaya\~{}en, yaya\~{}en'' vang lên trầm ấm, đủ khiến bất kỳ đứa trẻ nào ngủ quên giữa ban ngày. Nhưng đôi mắt đen láy, trong veo ấy của em bé lại nhìn thẳng vào tôi, như thể em vừa phát hiện ra một hạt bụi thời gian lạc lõng. Em mỉm cười, không hề biết rằng chỉ vài thập kỷ nữa thôi, chính mảnh đất này sẽ chứng kiến bao biến động, và nụ cười của em sẽ hóa thành tiếng khóc vọng lại giữa rừng thông.

Gió thổi nhẹ, cánh hoa mơ rơi xuống đầu tôi. Tôi đỡ lấy, nắm chặt trong lòng bàn tay - như giữ lấy mảnh ký ức cuối cùng còn nguyên vẹn, trước khi bước tiếp vào màn đêm đã định sẵn.

Tiếng ầm ì vang lên từ sâu thẳm lòng đất, như có con thú ngủ quên vừa giật mình trở mình. Từng tiếng rầm rầm rung nhẹ dưới đế bàn chân, rồi bỗng chốc nhanh hơn, mạnh đến mức cánh hoa mơ còn vương trên mái hiên rụng lả tả. Cả làng như bị ai giật mất nắp, tiếng người vỡ òa: ``\textbf{Động đất!}''

Yuko-sensei hét lên trước tiên, nhưng giọng cô cũng chỉ là giọt nước rơi vào chảo dầu đang sôi. Suzune-san cũng đang hoảng loạn không kém, khi cô nàng hoảng loạn chạy lòng vòng, chỉ để rồi bị Sakura-san giữ lại: ``Bình tĩnh nào, Karamomo-san!'' Dẫu biết rằng chúng tôi sẽ không bị ảnh hưởng bởi sự hỗn loạn trước mắt, nhưng nhìn cảnh tượng mọi người chạy tán loạn và những tiếng la thất thanh, tôi không thể không cảm thấy sởn gai ốc trước sức mạnh của Mẹ Thiên nhiên. Người lớn tháo chạy, trẻ con khóc thét, đàn bà gọi chồng gọi con bằng những âm thanh cắt đứt không khí. Một mái rạ ở cuối xóm gãy rạc, rơi xuống như lá bài được lật ngửa. Từ triền đất sau lũy tre, khối đất nâu sẫm trượt xuống, cuốn theo lối đi nhỏ và cả bóng dáng của những người đang bỏ chạy.

Kaigou-chan nghiêng đầu, mắt vẫn dán vào dòng người: ``Không chỉ động đất\dots\ mà còn sạt lở nữa.'' Cô nàng VU nữ gật nhẹ, giọng khàn đặc: ``Nhìn cách họ xoay xở là biết, tai họa này đã quay lại nhiều lần, đủ để họ nhớ kỹ từng bước chạy.''

Trong mắt họ không chỉ hoảng loạn, mà còn một thứ gì đen đặc hơn: nỗi sợ đã ngấm vào xương. Tiếng kinh cầu vang lên rời rạc, như những mảnh vải trắng được giật ra từ bức tranh cũ, không dính kết, chỉ bay lơ lửng giữa bụi đất.

Và rồi, giữa tiếng người, tiếng gió, tiếng đất gầm, một niềm tin lạ lẫng lách qua khe hở: ``Có thứ gì đó đang trách mắng chúng ta,'' tôi thì thầm. Yuko-sensei không nhìn tôi, chỉ nhìn theo bà cụ quỳ giữa đường, tay bưng mâm cơm cúng đất: ``Khi lý trí cạn kiệt, người ta dựng thần linh từ chính nỗi sợ của mình.''

Khung cảnh đột ngột thay đổi. Ánh sáng ban ngày như bị ai vặn nút, tắt ngóm giữa chừng; bầu trời không trăng, không sao, chỉ còn lớp màn đen dày đặc ép xuống nặng trĩu. Gió cũng ngừng thở, lá thôi xào xạc, như thể cả khu rừng đồng loạt bịt miệng chờ điều gì đó sắp lộ diện. Chúng tôi đứng giữa một khoảng trống tròn trĩa, quanh quẩn bởi những thân cây cao lênh khênh, như khán giả vô hình đã đến đông đủ để chứng kiến một vở kịch cổ.

Giữa khoảng trống ấy, đất đã được xới tung, tạo thành một hố nông chừng nửa mét. Thành đất tươi rói, còn vương mùi tanh của gốc rêu vừa bị cắt đứt. Đáy hố tối như miệng cống bị bịt kín bằng lớp vải đen; nhìn kỹ mới thấy những dải vải trắng buộc hờ trên miệng hố, như chỉ cần một cơn gió nhẹ là sẽ tuột xuống, phủ kín điều gì đang nằm bên dưới.

Từ phía bìa rừng, những bóng người lớn bắt đầu lần lượt xuất hiện. Họ mặc áo nâu sẫm, tay cầm đèn dầu, bước đi nặng nề như mang theo cả núi đất trên vai. Giữa họ, từng đứa trẻ được dắt đi, không dắt bằng tay mà bằng dải lụa trắng quấn cổ tay, như những con thú nhỏ được dắt đi dâng lễ. Trẻ con không khóc, không la hét, chỉ mở to đôi mắt đen nhánh nhìn quanh, như thể cả thế giới vừa bị ai bật nút ``tắt tiếng''. Một bé trai ngã về phía trước, lập tức bị người lớn kéo lại, không dữ dằn, chỉ nhanh gọn như sợ đứt mất dây cương vô hình.

Suzuki nắm chặt tay tôi, ngón tay lạnh toát. Em thì thầm, giọng run như sợ câu nói sẽ rơi xuống hố: ``Trông các em ấy không ai tự nguyện cả\dots\ phải không?''

Yuko-sensei đứng gần mép hố, hai tay siết chặt thành nắm đấm. Cô nói từng tiếng, giọng đều đến mức khiến không khí rung nhẹ: ``Đây không phải nghi thức. Đây là một vụ giết người được trang điểm bằng lời cầu kinh.''

Shion-san đứng phía sau, mặt cúi thấp, bóng áo choàng đen phủ lên mặt đất như tấm khăn tang khổng lồ. Cô cất giọng trầm ấm, lần đầu tiên không còn chút đùa cợt: ``Nhưng với họ, đây là cứu rỗi. Họ tin rằng chỉ cần chôn sống một phần, còn lại sẽ được tha.''

Mari-san đứng cách đó hai bước, mắt không chớp, nhìn vào hố như nhìn vào một chiếc gương vỡ. Cô lắc đầu chậm rãi, từng cái một, như thể mỗi lần lắc là một lời từ chối vô thanh. ``Không\dots\ không phải như thế này\dots'' Cô nói nhỏ đến mức chỉ có môi cử động, như sợ tiếng nói sẽ làm vỡ thêm điều gì đang treo lơ lửng trên đầu chúng tôi.

Trong hàng ngũ những đứa trẻ bị dắt đi, một bóng nhỏ bước chậm rãi ra khỏi bóng tối. Không cần ai chỉ tên, tôi đã nhận ra, đó là \emph{cô bé ấy }. Mái tóc đen tết đuôi sam lúc lắc theo từng bước chân, đôi mắt xanh lục trong veo như hồ ngọt giữa rừng sâu, nhưng lần này không còn lấp lánh ánh trăng nữa. Ánh nhìn ấy mờ đi, lặng lẽ, như thể đã nhìn thấy trước cả kết cục của chính mình.

Em không khóc. Không giật mình. Không cố chạy trốn. Em chỉ đứng đó, giữa vòng vây những người lớn, như một bức tượng nhỏ bằng xương và gió.

Tôi không thể rời mắt khỏi cô bé mắt xanh ấy. Không phải vì tôi kinh hãi gì em ấy cả. Mà vì trong đôi mắt long lanh tưởng như ngây thơ ấy, không hề có một đứa trẻ bình thường. Chỉ có một ký ức đang nhìn lại tôi.

``C-Cô bé đó\dots'' Suzune-san thì thầm, giọng nhỏ như tiếng lá rơi chạm đất. ``Là em ấy\dots\ phải không?'' Suzuki siết chặt tay tôi, móng tay gần như cắt vào da.

Những người lớn bắt đầu di chuyển. Họ xếp thành vòng tròn quanh hố. Một ông già cầm chiếc roi mây khẽ gõ xuống đất—một nhịp.

Hai nhịp.

Tiếng gõ như mở màn cho một vở kịch đã diễn đi diễn lại hàng trăm năm. Đất được xúc lên, từng nắm, từng nắm một, rơi xuống như mưa. Không ai nói gì, chỉ có tiếng đất va vào nhau, và tiếng hơi thở của những đứa trẻ đang cố giữ lấy không khí cuối cùng.

Trong tiếng cuốc xẻng và tiếng gõ gậy đó, em vẫn không nhúc nhích. Chỉ nhìn. Nhìn vào hố. Nhìn vào họ.  Nhìn vào chúng tôi, như thể em đang khắc gương mặt mỗi người vào một nơi sâu thẳm không thể xóa nhòa.

``Tại sao em ấy không chạy?'' Azuki-san lẩm bẩm, giọng rung như dây đàn chùng.

``Vì em đã biết,'' Shino-san đáp, mắt không rời khỏi đôi mắt xanh lục kia, ``không ai sẽ cứu em ấy.''

Ánh nhìn ấy\dots\ không thể hiện sự oán hận. Càng không có sự cầu xin nào. Chỉ ghi nhớ. Như thể em đang thu thập từng khoảnh khắc này, để mai sau, khi đất đã lấp kín, khi tiếng khóc đã tắt, khi người ta đã quên, rằng em vẫn còn đó.

Sakura-san lùi lại một bước, môi mấp máy: ``Chính khoảnh khắc này\dots\ là mốc thời gian. Là nơi mọi thứ, từ lời nguyền trong trường, cho tới những vòng lặp kia\dots\ đã bắt đầu.''

Khung cảnh như bị ai bấm nút tạm dừng. Gió ngừng lùa qua tán thông, tiếng lá khô im bặt, chỉ còn mùi đất xới bay lên nồng nặc. Tôi nghe rõ nhịp tim mình dồn dập, xen lẫn tiếng thở gấp của Suzune sau lưng. Tất cả đang chờ một lời giải thích cho cơn ác mộng vừa hiện ra.

Tôi lên tiếng trước, giọng khô khốc như vừa nuốt phải bụi đất: ``Chuyện này\dots\ không phải trùng hợp. Đây chính là gốc rễ của mọi thứ đang trói buộc chúng ta, không chỉ cả ngôi trường, mà có lẽ\dots\ là toàn bộ thị trấn này.''

Kaigou-chan đồng tình, mắt vẫn dán vào hố đen nơi đám trẻ vừa bị lấp. Cô nàng đưa tay lên miệng như muốn nén cơn buồn nôn: ``Một lời nguyền không tự nhiên mà có. Con người đã tạo ra nó, bằng chính đôi tay mình.''

Suzuki siết nhẹ tay tôi, đầu ngón tay lạnh toát. Giọng em nhỏ tới độ gần như van xin: ``Họ\dots\ họ đã cố hy sinh những đứa trẻ để tránh thiên tai phải không?''

Mari-san bước lên, giày cô nặng trĩu bùn đất. Cô nhìn quanh như đang đọc một bản nhạc thiếu nốt, rồi khẽ nói: ``Nhưng đây là điều khiến tôi không hiểu. Nếu nghi thức thành công, động đất và sạt lở phải dừng lại ngay sau đó rồi chứ nhỉ?''

Một khoảng lặng ngắn, đủ để nghe tiếng cô bé kia thở ra sau lớp sương mù. Tôi tiếp lời, lưỡi như dính vào vòm miệng: ``Thì, đáng lý là thế. Cơ mà, trong biên bản lưu trữ mà Shion-san cầm, họ vẫn ghi nhận những vụ sạt lở xảy ra trong những năm tiếp theo. Dân làng phải di dời sau đó.''

Yuko-sensei khoanh tay lại, khẽ gật: ``Đúng vậy. Nghi thức không những vô ích, mà còn để lại di chứng.''

Kaigou-chan vung tay, như muốn xua tan mùi máu tanh còn vương trong không khí: ``Vậy nghĩa là họ đã thất bại hoàn toàn. Không chỉ không cứu được làng, họ còn tạo ra một cái bẫy tinh thần.''

Sakura-san lắc đầu, mái tóc cô dính đầy sương đêm: ``Không đâu. Nó còn tệ hơn thế. Họ đã tạo ra một ``điểm neo''. Một chỗ dừng chân không lối thoát cho linh hồn những đứa trẻ.''

Tôi quay sang nhìn cô bé kia. Em vẫn đứng im, mắt xanh lục như hai hồ nước tối. Tôi cảm thấy lạnh sống lưng sau kết luận có phần thẳng thắn và khó nghe ấy. Những linh hồn trẻ con ấy đã bị giữ lại, và không được siêu thoát, tất cả chỉ vì sự mù quáng và ích kỷ của người lớn.

Shion-san bước lên, áo choàng đen phần phật như cánh dơi vừa thức giấc, đôi mắt xanh-nâu lóe ra dòng chữ cổ trên không trung chỉ cô mới thấy: ``Trong mười ba thế kỷ bị phong ấn, oán khí của các `hồn thú nhỏ' đã ngưng tụ thành `Hồ Bùn Đen Vô Tận,' một chấm đen trong Ma trận Thời gian. Chúng không tan, không phai, chỉ đặc quánh thêm\dots\ rồi biến thành `Vòng Lặp Định Mệnh' mà chúng ta đang bị mắc kẹt! Ta nói thế có đúng không?''

Tôi nhớ lại lần đầu tiên bước vào vòng lặp, khi lời đồn về việc nếu không ``đối xử đúng'' với người mới, họ sẽ bị ``gỡ bỏ'' khỏi thế giới này, như xóa một dòng kẻ bảng. Ban đầu, đó chỉ là một lời đồn đại vô thưởng vô phạt nhằm dọa những kẻ hú vía, chỉ để rồi mọi thứ đều sụp đổ trước mắt, và cũng từ đó biết được quá khứ của Misora Shino-san như thế nào.

Rồi những người khác cũng biến mất khi họ đi qua nhà ga Aogami và không bao giờ trở lại. Tiếng tàu rít lúc nửa đêm không phải tàu chở khách, mà là tàu chở những linh hồn bị lãng quên, chạy thẳng một chiều về một nơi hư không, nơi những kẻ bị lạc lõng được chào đón.

Và ``thế giới hạnh phúc'' giả tạo ấy, tôi đã sống trong nó bao nhiêu lần? Một buổi sáng nắng nhẹ, lũ bạn cười đùa trên hành lang, kết bạn với nhiều bạn học, tất cả đều hoàn hảo đến nghẹt thở. Nhưng khi tôi bước ra khỏi vòng lặp, tôi nhận ra: cái gọi là ``hạnh phúc'' ấy chỉ là lớp vỏ bọc mỏng manh của một lời nguyền không muốn ta nhớ đến nỗi đau. Mỗi lần tôi cố gắng kể cho ai nghe về những người mất tích, thế giới lập tức ``reset'', quay ngược về ngày đầu tiên của học kỳ, khi tất cả còn nguyên vẹn, ngoại trừ ký ức của tôi. Và tôi hiểu: tất cả những vòng lặp ấy, tất cả những lần mà Sora muốn thực hiện mong ước nhỏ nhoi này, mọi thứ đều bị sụp đổ và quay về con số không.

Giờ đây, tôi đã biết những hiện tượng trên đều bắt đầu từ nghi thức mù quáng này, từ hố đất nhỏ nơi khu rừng đã từng nuốt chửng những đứa trẻ vô tội để đổi lấy một lời hứa hão huyền. Khu rừng không chỉ giữ xác họ; nó giữ cả thời gian, giữ cả những kẻ còn sống trong vòng xoáy lặp đi lặp lại, để chúng tôi mãi mãi không thoát khỏi nỗi sợ ``nếu không đối xử tốt, người kia sẽ biến mất''. Bởi vì chúng tôi - tựa như những đứa trẻ bị bỏ lại - chính là phần ``còn lại'' của lời nguyền.

Suzuki thì thầm, nước mắt lăn dài trên má: ``Mười ba thế kỷ\dots\ bao nhiêu thế hệ nữa sinh ra và chết đi, còn các em vẫn mắc kẹt nơi này.''

Tôi cảm thấy cổ họng mình nghẹn đắng: ``Và rồi, khi lớp vỏ kiếp người không còn chứa nổi, nó vỡ ra thành thứ chúng ta vẫn gọi là `lời nguyền.'''

Không ai trong chúng tôi phản đối, bởi những gì chúng tôi thấy trong miền ký ức méo mó kia\dots\ là thật. Chỉ có gió thở ra một hơi dài, như thanh minh rằng nó chỉ là kẻ chứng kiến vô tội.

Tiếng lá khô lại rạo rực, như có ai vừa bật công tắc âm thanh. Không gian rung lên, không phải rung của gió mà rung của ánh sáng, một luồng trắng nhạt lan từ tâm hố, làm méo mó cả bóng cây. Những mảnh ký ức vừa được gọi dậy bắt đầu nứt vỡ, từng mảnh như thủy tinh rơi xuống lòng đất, lóe lên rồi tắt ngóm.

Giữa tàn ảnh tan rã, cô bé hiện ra lần cuối. Làn sương đen cuộn quanh em như tấm màn nhung đã cũ, chỉ chừa lại khuôn mặt nhỏ nhắn, đôi mắt xanh lục nhìn thẳng vào chúng tôi, không còn u uất, chỉ còn mỏi mệt. Em mấp máy, giọng mỏng như lá khô vừa rụng:

``Mọi người\dots\ đã thấy hết rồi, phải không?''

Tôi chưa kịp gật, chưa kịp nói một tiếng. Không gian như tấm gương bị búa tạ đập trúng, vỡ vụn thành vô số thời khắc. Tiếng em, tiếng gió, tiếng tim tôi, tất cả đứt quãng, rơi xuống hố sâu không đáy.

Tôi chỉ kịp nhắm mắt.

Khi tôi mở mắt lại, trời đã trở lại màu đen rừng thông. Chúng tôi đứng chật cứng giữa vòng trũng quen thuộc, lòng bàn chân vẫn cảm giác mòn mỏi của lớp lá mục. Không còn khói bếp, không còn tiếng chuông gió, chỉ còn mùi đất tươi xới lên, và cái nhìn của cô bé in sâu vào đêm, như vết sáng hằn trên mặt tường đã đổ.

Im lặng bao trùm, dày đặc đến nỗi tiếng thở của chính mình cũng như vang lên từng hồi trong đầu. Tôi cúi nhìn bàn tay, ngón tay vẫn rung nhẹ như còn giữ mãi hình ảnh những dải lụa trắng quấn cổ tay. Mùi đất tươi vừa xới vẫn bám trên da, nồng nặc đến khó chịu.

Azuki-san lên tiếng trước, giọng khàn như thể cát vừa rơi vào cổ họng: ``\dots Vậy là\dots\ cái này, không phải kiểu lời nguyền `tự nhiên' gì đó nhỉ.''

``Đúng là như vậy,'' tôi đáp. ``Nó có đầu đuôi rõ ràng.''

Kaigou khoanh trước ngực, mắt nhìn xa xăm về phía đồi thông: ``Họ tự tay dựng đàn, tự tay chọn con chiên. Và cũng chính họ phải lãnh hậu quả sau này.''

Sakura-san bước lên, mũi giày chạm mép hố, cô khẽ nói như sợ đất dưới chân sẽ nứt ra: ``Cái chết không đáng sợ bằng cái chết bị gọi là `lễ vật'.''

Tôi tiếp lời, mỗi chữ như nhỏ giọt vào không khí lạnh: ``Những đứa trẻ ấy bị chọn. Chúng bị đưa đi, không có quyền lên tiếng, không được thanh minh, và rồi\dots''

``\dots bị những bàn tay mu muội ấy chôn sống khi vẫn chưa hiểu nổi tại sao bản thân lại rơi vào tình cản như vậy,'' Shino-san tiếp.

Yuko-sensei siết chặt tay áo, khóe môi run nhẹ: ``Dưới danh nghĩa cứu rỗi, họ chỉ cứu lấy nỗi sợ của chính mình.''

Suzuki kéo nhẹ tay tôi, giọng em nhỏ đến mức gió thổi qua cũng có thể cuốn mất: ``Không ai dám thốt ra hai tiếng `xin lỗi'\dots\ để lũ trẻ kia có chỗ về. Cũng không một ai thừa nhận chính tay họ đã giết chết những đứa trẻ ấy.''

Shino-san cúi đầu, mái tóc che kín nửa khuôn mặt, cô thì thầm như đọc một câu thần chú đã mòn vẹt: ``Chính vì không được nhìn nhận, nên họ không được buông tha, và vẫn ám ảnh nơi này đến tận bây giờ.''

Tôi quay đầu lại. Trong màn đêm đặc quánh, những chấm sáng li ti bắt đầu nhú lên như hạt sương bị lửa nung. Chúng run rẩy, ẩn ẩn hiện hiện, không còn trốn sau tán thông nữa. Chúng lơ lửng gần chúng tôi, gần đến mức tôi nghe được tiếng thì thầm: ``Chúng tôi đang ở đây''.

Không phải oán hận. Chỉ là một lời xin được nhìn thấy, được gọi tên, lần đầu tiên sau mười ba thế kỷ bị lãng quên.

Shion-san đứng yên, lần đầu hỏi một cách nhã nhặn, tạm rời xa cái tính ảo tưởng của mình: ``Vậy cái lõi rốt cuộc là gì?''

Tôi thấy mình trả lời trước khi kịp nghĩ. Chỉ bốn từ ngắn gọn: ``Là sự phủ nhận.''

Tôi ngừng một nhịp, để tiếng mình không vỡ, và rồi tiếp tục: ``Người lớn không bao giờ nói: `Chúng tôi đã giết các em'. Họ gọi đó là lễ, là phong tục, là cách xin trời tha. Nhưng không ai đứng ra thừa nhận: những đứa trẻ ấy đã chết oan. Không ai viết nên hai chữ `kết thúc' cho chúng, mà đúng hơn, chúng không xứng đáng bị đối xử như vậy.''

Sakura-san gật chậm, như gõ nhịp cho lời tôi. ``Nên chúng không được siêu thoát, thậm chí còn kéo dài dai dẳng tới bây giờ.''

Kaigou-chan khẽ gật, mắt vẫn dán vào những đốm sáng lơ lửng như đom đóm trắng.
``Vậy nên\dots\ chúng chỉ cần lời thú nhận?''

Tôi thở ra, hơi nóng bốc lên giữa đêm lạnh. ``Không chỉ thú nhận. Chúng cần được ghi nhận rằng chúng đã từng tồn tại, đã từng bị hại, đã từng\dots\ là người.''

Suzuki nắm tay của mình chặt hơn, ngón tay run như lá thu. ``Em hiểu rồi. Nếu chúng ta nói tên chúng, kể câu chuyện của chúng, biết được sự tồn tại của chúng\dots\ đó có thể là câu trả lời. Và để làm được điều đó\dots''

``\dots chúng cần có những người lắng nghe.'' Yuko-sensei kết luận.

Shion-san ngẩng lần đầu, mắt ánh lên vẻ tỉnh táo lạ thường. ``Ký ức là then chốt. Nhưng phải là ký ức không thốt ra bằng miệng người lớn, mà bằng tiếng nói của chính những đứa trẻ.''

Không khí như ngừng đọng, nhưng lần này không phải là bầu không khí nặng nề của sự tuyệt vọng hay bế tắc. Mọi thứ đang lắng lại, như những mảnh vỡ của một bức tranh cổ đang tự tìm đường trở về đúng vị trí của nó. Ánh sáng từ đèn pin rung nhẹ trên tay của cô nàng Phù thủy tự xưng, bóng của cô loang ra trên mặt đất như một vệt mực cũ đang tìm lại dòng chữ đã mất.

``Khoan đã,'' cô bất ngờ lên tiếng, giọng cô nhỏ nhưng đủ khiến tất cả quay lại. Cô lật vội cuốn tài liệu đã bị gập lại từ lúc nãy, những trang giấy mỏng manh rung lên trong gió đêm. ``Hình như trong đây có điều gì đó rất quan trọng.''

Cô dừng lại, ngón tay dừng giữa chừng, rồi chỉ thẳng vào một dòng chữ nhỏ, gần như bị lẫn trong lớp bụi thời gian trong dòng công thức cô xưa: ``Nghe này, `Người được chọn phải tự bước vào.'''

Azuki-san nhướng mày, mắt cô nheo lại như đang cố nhìn xuyên qua bóng tối để tìm ý nghĩa ẩn sau dòng chữ ấy. ``Tự bước vào? Ý họ là\dots\ tự nguyện?''

Shion-san gật đầu, mắt cô không rời khỏi trang giấy. ``Không có sự ép buộc. Không có lệnh truyền. Nếu đứa trẻ không tự đặt chân vào hố, hoặc miễn cưỡng vào đó, thì toàn bộ nghi thức sẽ vô nghĩa.''

Một làn gió lạnh thổi qua, như thể chính khu rừng cũng đang thì thầm lời xác nhận. Không phải họ đã giết những đứa trẻ, mà họ đã khiến chúng tự bước vào cái chết. Không bằng bạo lực, mà bằng lời ru, bằng lời hứa hão huyền, bằng niềm tin rằng máu của mình sẽ cứu lấy làng, rằng nước mắt của mình sẽ thành mưa, rằng sự hi sinh của mình sẽ được trời ghi nhận.

``Vậy thì\dots\ đây không chỉ là một lời nguyền bình thường đâu,'' Shino-san thì thầm, ``đây là một lời dối trá được đắp bằng lòng tin một cách đầy trơ trẽn.''

Shion-san khép cuốn tài liệu lại, mắt cô nhìn xa xăm về phía hố đất đã lấp kín từ lâu. ``Và nếu vậy\dots\ thì chúng ta đang không chỉ đối mặt với oán khí của những đứa trẻ. Chúng ta đang đối mặt với một lời hứa đã bị phản bội, và không ai dám thừa nhận điều đó.''

Azuki-san chậm rãi hít một hơi, mắt vẫn dán vào hố đất tối om trước mặt. Gió đêm lùa qua kẽ lá, làm rung nhẹ tấm vải trắng còn vương trên miệng hố, như thể chỉ cần một cú giật tay là cả ký ức sẽ lộ ra. Cô lên tiếng, giọng khàn vì bụi đất: ``Vậy\dots\ chúng ta phải làm lại nghi thức?''

Tôi lắc đầu, đôi mắt vẫn không rời khỏi chỗ đất vừa bị xới tung. Không khí quanh đây nặng như bị ai đổ chì vào gió. ``Cũng gần như vậy, nhưng không phải là để hiến tế,'' tôi sửa lại, ``mà là `hoàn tất'. Đưa họ về nơi họ thuộc về. Để họ được đi.''

Shino-san gật nhẹ, mái tóc dài đen nhánh che khuất nửa khuôn mặt. Cô thì thầm như đang nói với chính mình: ``Đưa họ rời đi. Không phải đẩy họ xuống, mà là dắt họ lên.''

Một khoảng lặng ngắn. Suzune-san bước lên một bước, giọng run như khan sợ câu nói sẽ rơi vỡ: ``Nhưng\dots\ ai sẽ làm điều đó? Mình không biết đầu đuôi ra sa, nhưng có lẽ\dots\ không phải ai cũng\dots\ có thể, đúng không?'' Cô nhìn quanh, mắt dừng lại trên từng khuôn mặt. Ánh nhìn không cầu xin, chỉ là một lời xác nhận.

Tôi gật. Giọng mình trầm xuống, như đang nói với cả khu rừng đang nghe lén: ``Đúng. Không phải ai cũng phù hợp. Chỉ những người đã từng bị dòng đời này níu lại. Những người đã từng\dots\ liên kết với hiện tượng này.

Sakura-san tiếp lời, mắt cô nhìn xa về phía đồi thông, nơi những dải vải trắng vẫn còn phấp phới trong ký ức: ``Những người đã từng bị lời nguyền chạm vào da thịt. Những người đã từng\dots\ mất một phần mình ở đây.

Shino-san cúi đầu, giọng nhỏ đến mức gần như tan trong không khí: ``\dots Là mình.

Cô nàng Vu nữ quay lại, mắt cô không chớp: ``Tôi cũng vậy.

Kaigou liếc sang tôi, nở một nụ cười hơi mếu như biết được số phận. Cô khẽ nói: ``Vậy là chúng ta không thoát khỏi vai này rồi.

Tôi nhìn sang Suzuki. Em đứng im, hai tay nắm chặt vạt áo. Ánh mắt em không sợ hãi như thể em đã biết từ trước.

``Tôi, cô, và em ấy, tôi nói, ``ba người chúng ta. Và hai người họ.''

Suzuki gật nhẹ, giọng em nhỏ như tiếng lá rơi: ``Vâng.''

Azuki-san đếm nhanh bằng đầu ngón tay, mắt cô lia qua từng người: ``Một, hai, ba\dots\ Vậy là năm người rồi.''

Cô quay sang nhóm còn lại, giọng hơi cao: ``Còn bọn mình thì sao?''

Yuko lên tiếng trước tiên, giọng cô cứng rắn như đang đóng cửa một căn phòng:

``Chúng ta sẽ đứng ngoài. Các em không đủ điều kiện.

Shion-san chép miệng, đôi mắt xanh-nâu nhìn xuống cuốn tài liệu trong tay như thể nó vừa phản bội cô: ``Và ta cũng vậy. Chán thật đấy.''

Kaigou liếc sang cô, môi cô nhếch lên một nụ cười không vui: ``Tưởng điên điên thế nào mà cũng hữu dụng đấy chứ nhỉ?''

``\dots Ờm\dots\ Cảm ơn?'' cô nàng nhíu mày, không rõ là đang trêu hay nói thật.

Gió đêm bỗng nhẹ hẳn, không còn se lạnh như trước, mà như một hơi thở vừa được thả lỏng sau bao năm nín ké. Tôi ngửi thấy mùi cỏ lá hòa vào đất ẩm, thoang thoảng một chút hương thông già, như thể khu rừng cũng đang thở ra nhẹ nhõm. Tôi lùi một bước, để lưng tựa vào thân cây tùng thấp, mắt vẫn không rời bóng cô bé đứng giữa vòng tròn tối om.

``Vậy là rõ rồi,'' tôi thì thầm, giọng khàn vì bụi đêm. ``Không cần ai phải chết nữa.''

Tôi ngừng lại, để câu nói bay theo làn khói nhỏ giữa hai hàng mi. ``Chỉ cần… làm đúng thứ lẽ ra đã xảy ra từ đầu.''

Cơn gió lạnh lướt qua vai tôi, mang theo hơi ấm lạ lùng của ánh đèn pin vừa tắt. Kỳ lạ thay, tôi không cảm thấy lạnh. Tôi ngẩng lên. Cô bé vẫn đứng đó, mái tóc đuôi sam hơi xõa, đuôi tóc đung đưa như cái chớp mắt chậm rãi. Đôi mắt xanh lục không còn đặt ở chân trời xa, mà đang ở ngay trước mặt chúng tôi: gần hơn, nhưng vẫn giữ một khoảng cách mỏng, đủ để cảm thấy còn e dè.

Tôi không nói thêm. Tôi chỉ nhìn lại, để mắt mình nói thay: chúng tôi không định bỏ đi. Một khoảng lặng tròn trịa lấp đầy khoảng trống giữa hai bên, như mặt hồ vừa ngừng gợn sóng.

Rồi cô bé khẽ mở miệng. Tiếng cô nhỏ đến mức tôi phải nín thở mới nghe được: ``Nếu\dots\ mọi người thật sự muốn giúp\dots''

Em ấy chợt dừng, mắt lướt qua từng khuôn mặt của những người được chọn - Shino-san đang cúi, Sakura-san đang siết chặt tay áo, Suzuki-san run nhẹ bên tôi, Kaigou-chan hiện rõ nét mặt quyết tâm qua cái gật đầu nhẹ - rồi dừng lại ở tôi, như gieo một hạt giống vào đáy mắt.

``\dots thì đừng quên chuyện này.''

Nghe không giống một câu lệnh, cũng không phải van xin, chỉ là lời nhắc nhở của một đứa trẻ đã quá quen với việc bị lãng quên, nhưng lần đầu tiên sau một khoảng thời gian rất dài đã bắt đầu nhen nhóm sự mong mỏi.

Tôi gật đầu, không cần nghĩ ngợi thêm. ``Bọn anh sẽ không quên đâu.''

Suzuki siết nhẹ tay tôi, mắt em vẫn dán vào cô bé kia. Em gật chậm, như thể mỗi chuyển động đều nặng hơn cả trái đất. ``Bọn chị nhớ rồi.''

Ánh sáng li ti bỗng chụm lại, tạo thành một vòng tròn trắng mỏng quanh hố đất. Không phải ánh sáng của đèn, cũng không phải lửa. Đó là ánh sáng của lời hứa vừa được thốt ra. Dù rất mong manh, nhưng đủ để xua tan bóng tối mười ba thế kỷ.

Cô bé nhìn chúng tôi thêm giây lát, rồi bước lùi một bước, nhường chỗ cho điều gì đó sắp tới. Sương đen cuộn lại, che khuất em từng chút, từng chút, để lại chỉ còn đôi mắt xanh lục lấp lánh giữa màn đêm, như ngọn đèn cuối cùng trước khi tắt.

Gió lại thở. Lá khô xào xạc. Không ai nói thêm lời nào nữa. Cánh rừng Aogami đã quay trở về với nét hoang sơ vốn có, như chưa hề có một bóng ma nào xuất hiện.

Không ai trong chúng tôi lên tiếng, trong số chúng tôi vẫn có những người còn hơi sốc, như thể những gì vừa xảy ra chỉ là trong cơn mơ sảng; chỉ tiếc rằng đây là sự thật. Sự im lặng của những người vừa cùng nhau nuốt xuống một khối u nặng. Tôi bước lên triền đất cao, đá lở dưới gót giày rơi xuống vùng trũng như những mảnh ký ức không kịp mang theo. Không ai ngoái lại, nhưng tôi biết, cả nhóm đều đang nghe tiếng vọng của chính mình vang lên từ đáy hố vừa bỏ lại.

\ct

Suzune-san lẽo đẽo cuối đoàn, vai co lại, bước chân ngắn cũn như đang đi trên mặt phẳng không ổn định. Cô vốn hay ríu rít, giờ chỉ còn tiếng lá khô dưới đếi giày. Đèn pin trên tay cô rung lên, quét vài vệt sáng loãng lên thân cây rồi tắt ngúm. ``\dots Mình từng tưởng chỉ là cuộc dã ngoại đêm khuya\dots'' giọng cô vẫn còn run run, ``\dots ngờ đâu lại thành một chuyến đi đáng sợ\dots\ như vậy\dots''

Mari-san đi sát bên, ôm cuốn sketch như cái phao. Trang giấy cuối cùng cô mở ra vẫn còn in nét phấn trắng của bãi đất lúc nãy; cô không lật tiếp, chỉ siết sổ vào lòng, mắt dán xuống đường mòn. ``Giờ tôi mới hiểu,'' cô khẽ nói, ``tại sao tôi lại vô thức vẽ nên bức tranh đó.''

Phía trước, Azuki-san dừng bước, hai tay xoa mặt như gỡ bỏ lớp mặt nạ bụi than đêm. Cô thở dài, hơi nóng bốc lên giữa khuya: ``\dots Ừm. Vậy là\dots\ không phải kiểu phiêu lưu vui vẻ nhỉ.''

Shion-san đi ngang, tập tài liệu gõ nhẹ vào đầu gối với mỗi bước xuống dốc. Giọng cô thấp hơn bình thường, không còn vẻ ``bật đèn pha'' như mọi khi. ``Ngươi tưởng rừng đêm là chỗ picnic?'' cô khẽ cười cợt, nhưng cái cười chỉ nửa miệng. ``Nhưng mà nha,'' cô dừng lại, quay đầu nhìn về phía vùng trũng đã khuất sau lũy thông, ``chúng ta vẫn chưa thoát khỏi bóng ma đó đâu. Trận chiến mới chỉ bắt đầu thôi.''

Rồi như ai bấm nút, cô như bật lại chế độ ``nghiên cứu'': ``Ta về căn cứ! Lục toàn bộ thư tịch cổ! Giải mã nghi thức!''

Yuko-sensei bước tới, khoác vai cô học trò đó, giọng đều đều không mấy ngạc nhiên: ``Cụ thể là em sẽ lục mớ giấy trong phòng CLB, phải không?''

Cô nàng Phù thủy hất tay, mặt vênh lên nhưng không đủ sức: ``\dots Đúng.''

``Vậy còn chờ gì nữa.'' Đôi mắt của cô chủ nhiệm dõi xa xăm về phía rừng thông, không còn vẻ nghi hoặc như trước, chỉ là một tĩnh lặng chấp nhận: cô đã đứng trong câu chuyện này, dù không phải diễn viên chính.

Tôi bước xuống dốc, cát khô lún nhẹ dưới đế giày. Kaigou-chan đi bên trái, tay cô đút xuống túi váy đồng phục, mắt nhìn thẳng vào màn đêm đen như nỗi đang vật lộn trong lòng. Suzuki khoác tay tôi bên phải, bước chân nhỏ hơn, nhưng vẫn cố giữ nhịp. Ba đứa tạo thành một hàng ngang chật hẹp, tiếng lá thông khô nứt dưới chân là âm thanh duy nhất.

Chúng tôi không trao đổi lời nào. Mà cũng chẳng cần. Gió đêm thổi qua mang theo mùi đất vừa xới, mùi ấy đã nói thay tất cả: chúng tôi vừa cùng nhau nuốt xuống một khối ký ức còn nóng hổi.

Cuối cùng, Kaigou-chan dừng lại trên bãi đất trống trước khi rẽ vào lối mòn về trường. Cô không quay sang, chỉ lên tiếng, giọng khàn vì bụi đêm: ``Ít ra cũng biết được tại sao mọi chuyện lại thành thế này.''

Tôi khẽ chẹp miệng, mắt vẫn dán vào đường chân trời lờ mờ sau rừng thông. ``Ừ. Khởi đầu không phải do trời trừng phạt, mà do người lớn tự chọn lấy.''

Suzuki siết nhẹ tay áo tôi, ngón tay lạnh toát. Em ngẩng lên, đôi mắt long lanh trong bóng đèn pin đã tắt: ``Nhưng\dots\ làm sao để kết thúc thì\dots\ ta vẫn chưa rõ làm thế nào.''

Cả hai người lớn chúng tôi thở dài, đồng ý với ý của cô em gái đó: vấn đề mới chỉ được gạt được lớp nền, còn phần lõi thì chưa có cách nào tiếp cận được.

Phía sau chúng tôi, bóng hai người khác cũng vừa tới nơi bãi cỏ. Sakura-0san dìu Shino-san đi chậm, hai cái bóng dính sát vào nhau như hai mảnh ghép đang tìm cách trở lại hình cũ. Họ không lên tiếng, chỉ dừng lại cách chúng tôi vài bước, cùng nhìn về phía trường học ẩn trong màn sương đêm, nơi mà, sáng mai, chúng tôi sẽ phải trở lại với bộ mặt học sinh bình thường, dù trong lòng đã mang theo một hố chết vừa hé mở.

Chúng tôi bước qua bậc đá cuối cùng, để lại phía sau tán thông đen nhánh như bức màn đóng kịch. Ngay lập tức, làn gió đêm thô ráp thay bằng luồng không khí mát rượi của cánh đồng cỏ; hơi thở tôi không còn vướng mùi đất nữa, mà thoang thoảng hương cỏ non cắt dở. Cảm giác nặng nề như bị ai đó nhìn trộm cũng tan biến cùng tiếng cú kêu xa tít.

Cả nhóm tự động dừng bước, không ai bảo ai. Đèn pin tắt ngóm, chỉ còn ánh trăng khuya soi mỏng mảnh lên bả vai run rẩy của Suzune0san. Tôi nghe tiếng tim mình đập chậm lại, như thể cơ thể cũng nhận ra: đây là điểm dừng, chưa phải đích đến.

Azuki đập tay vào nhau, gượng cười: ``Vậy\dots\ mai gặp lại nhé!'' Giọng cô cố tỏ ra nhẹ nhàng, nhưng đuôi câu vẫn run nhè nhẹ.

Mari đứng cạnh, ôm cuốn sketch sát ngực. Cô gật, mắt vẫn dán vào bóng rừng sau lưng: ``Ừ, hẹn gặp mọi người.'' Dường như cô sợ nếu ngoái lại, bóng tối sẽ kéo cô trở vào.

Suzune khẽ vẫy tay, nở nụ cười hơi xệch, mồ hôi túa ra như vừa mới tắm: ``\dots Mình\dots\ sẽ suy nghĩ thêm.''

Shion-san giơ tập tài liệu lên, chất giọng tự tin: ``Cứ để Shion ta đây lo phần này! Trong mai mốt tin vui sẽ đến!''

Đoàn người rã dần, ai nấy lần lượt đi về một phía, tạm gác đi những gì vừa xảy ra trong khu rừng. Khi những tiếng bước chân dần lắng xuống trong màn đêm, chỉ còn năm bóng dáng đứng lại: tôi, Kaigou-chan, Suzuki, Sakura-san và Shino-san.

Không khí lúc này trở nên đặc quánh, như vừa được lọc qua lớp bụi thời gian. Sự im lặng bao trùm đến mức có thể nghe rõ tiếng lá khô rơi chạm nhẹ lên vai áo. Chúng tôi vô thức xếp thành một vòng tròn thưa thớt, mắt nhìn xuống mặt đất, ai cũng đang chờ đợi điều gì đó mà chưa ai dám lên tiếng trước.

Đúng vào khoảnh khắc ấy, cổ tay tôi khẽ rung lên.

*BÍP\dots\ BÍP\dots*

Tiếng đồng hồ vang vọng giòn tan và lạnh lẽo giữa màn đêm. ``Bước đầu đã hoàn thành.''

Giọng nói trầm ấm quen thuộc đột ngột cất lên từ chiếc vòng kim loại trên cổ tay tôi. Không còn là những tiếng bip khô khan nữa, mà là giọng nói của kẻ đã chủ động im lặng suốt từ đầu chuyến đi đến giờ. ``Tuy nhiên\dots\ dựa trên những gì các cậu vừa trao đổi, vẫn còn thiếu một điều kiện mấu chốt để thực hiện nghi thức giải thoát.''

Giọng nói ấy ngắn gọn, dứt khoát, nhưng mỗi từ đều mang sức nặng rõ rệt. Tôi cảm giác như cổ tay mình bị siết lại một chút, không phải vì chiếc vòng, mà vì chính nội dung lời nói.
``Điều kiện à\dots'' tôi lặp lại khẽ. ``Đúng là chúng tôi vẫn chưa tìm ra, nhưng\dots''

Shino-san lên tiếng, giọng rất nhẹ: ``\dots Nếu ký ức chính là chìa khóa\dots'' cô khẽ ngẩng đầu nhìn về phía khu rừng tối om, ``\dots thì chúng ta cần tái hiện lại nó.''

Sakura-san gật đầu nhẹ, giọng trầm tĩnh: ``\dots Một nghi thức,'' cô dừng lại một nhịp, như đang cân nhắc từng chữ, ``\dots nhưng không phải để hiến tế.''

Tôi hít sâu một hơi, mắt nhìn thẳng về phía trước, nơi bóng tối của rừng đang lan tỏa. ``Lần này\dots'' tôi nói, giọng kiên định hơn bao giờ hết, ``\dots chúng ta sẽ thực hiện được tới cùng.''

Sau những gì mà chúng tôi thấy trong màn gặp mặt đó, một ranh giới giữa cái đã qua và cái sắp tới, một khoảnh khắc mà chúng tôi đứng giữa hai thế giới, không thể quay lại, cũng chưa đủ can đảm để bước tiếp. Trong bóng tối, tiếng côn trùng kêu rả rích như đang đếm từng giây từng phút cho điều gì đó sắp xảy ra. Trong khoảnh khắc nàym sâu trong thâm tâm, chúng tôi đã biết rằng: từ thời điểm này, chúng tôi không còn là những con cừu ngơ ngác ngây thơ bị lạc hướng nữa, mà giờ chúng tôi đã tìm được hướng đi cho riêng mình\dots

\dots để kết thúc lời nguyền bao quanh thị trấn này, một lần và mãi mãi.

\end{document}